% Cours : Refaire l'impôt
% Compilation : pdflatex refaire-l-impot.tex (deux passes pour les renvois)
% Les figures sont les versions «~nues~» produites par  FIG_NU=1 python graphiques*.py
\documentclass[11pt,a4paper]{article}

\usepackage[T1]{fontenc}
\usepackage[utf8]{inputenc}
\usepackage{lmodern}
\usepackage[french]{babel}
\usepackage[a4paper,left=3.1cm,right=3.1cm,top=2.9cm,bottom=3.1cm]{geometry}
\usepackage{graphicx}
\usepackage[export]{adjustbox}
\usepackage{booktabs}
\usepackage{tabularx}
\usepackage{xcolor}
\usepackage{titlesec}
\usepackage[labelfont=bf,labelsep=period,font=small,singlelinecheck=false,
            justification=raggedright,skip=6pt]{caption}
\usepackage{fancyhdr}
\usepackage{enumitem}
\usepackage{microtype}
\usepackage{textcomp}
\providecommand{\euro}{\texteuro}
\usepackage[hidelinks]{hyperref}

\definecolor{encre}{HTML}{1A1A19}
\definecolor{bleu}{HTML}{003399}
\definecolor{vert}{HTML}{1F7A4D}
\definecolor{ocre}{HTML}{C79100}
\definecolor{brique}{HTML}{993333}
\definecolor{filet}{HTML}{C8C6BE}

\titleformat{\section}{\normalfont\large\bfseries\color{bleu}}{\thesection.}{0.6em}{}
\titleformat{\subsection}{\normalfont\normalsize\bfseries\color{bleu}}{}{0em}{}
\titlespacing*{\section}{0pt}{2.4ex plus 1ex minus .2ex}{1.4ex plus .2ex}
\titlespacing*{\subsection}{0pt}{2.2ex plus 1ex minus .2ex}{1ex plus .2ex}

\pagestyle{fancy}
\fancyhf{}
\renewcommand{\headrulewidth}{0.4pt}
\fancyhead[L]{\small\itshape Refaire l'impôt}
\fancyhead[R]{\small\thepage}

\setlength{\parindent}{0pt}
\setlength{\parskip}{0.72em}
\linespread{1.06}
\renewcommand{\arraystretch}{1.22}

% --- environnements maison -------------------------------------------------
\newcommand{\amorce}[1]{\textbf{#1}}

\newcommand{\partie}[2]{%
  \clearpage
  \vspace*{0.4cm}
  {\color{bleu}\hrule height 1.6pt}
  \vspace{10pt}
  {\footnotesize\color{ocre}\bfseries\MakeUppercase{#1}\par}
  \vspace{6pt}
  {\fontsize{19}{23}\selectfont\bfseries\color{bleu}#2\par}
  \vspace{8pt}
  {\color{bleu}\hrule height 0.6pt}
  \vspace{16pt}}

\newenvironment{encadre}[1]
  {\par\medskip\noindent\begin{minipage}{\linewidth}%
   \color{encre}\setlength{\parskip}{0.6em}%
   \hrule height 0.8pt width \linewidth \vspace{6pt}%
   {\footnotesize\bfseries\color{vert}\MakeUppercase{#1}}\par\vspace{-2pt}%
   \small}
  {\par\vspace{4pt}\hrule height 0.4pt width \linewidth\end{minipage}\par\medskip}

% L'argument d'un chapitre d'impôt, en une phrase, en tête de chapitre.
\newenvironment{argument}
  {\par\smallskip\noindent\begin{minipage}{\linewidth}%
   \color{encre}\hrule height 0.8pt width \linewidth \vspace{7pt}%
   \large\itshape}
  {\par\vspace{7pt}\hrule height 0.4pt width \linewidth\end{minipage}\par\medskip}

% Figure : légende au-dessus, image, puis sources et note de lecture en texte (fichier
% <figure>_note.tex écrit par les scripts de figures en mode FIG_NU).
\newcommand{\figbloc}[3]{%
  \begin{figure}[htbp]\centering
  \caption{#2}\label{#3}
  \includegraphics[max width=\linewidth,max height=0.42\textheight]{../figures/nu/#1.pdf}
  \IfFileExists{../figures/nu/#1_note.tex}{\par\vspace{4pt}%
    \parbox{\linewidth}{\raggedright\scriptsize\color{encre}\setlength{\parskip}{2pt}%
    \input{../figures/nu/#1_note.tex}}}{}%
  \end{figure}}

% ===========================================================================
\begin{document}
\addto\captionsfrench{\renewcommand{\figurename}{Figure}%
  \renewcommand{\tablename}{Tableau}}
\renewcommand{\figurename}{Figure}\renewcommand{\tablename}{Tableau}

\thispagestyle{empty}
\vspace*{1.2cm}
{\footnotesize\color{ocre}\bfseries COURS EN TREIZE CHAPITRES\par}
\vspace{10pt}
{\fontsize{30}{34}\selectfont\bfseries Refaire l'impôt\par}
\vspace{14pt}
{\large Pourquoi certains impôts coûtent beaucoup plus que d'autres pour le même euro de recette,
et ce que cela dit du système français.\par}
\vspace{12pt}
{\small\color{encre}Le cours est construit pour qu'on puisse réexpliquer chaque impôt à
quelqu'un d'autre. La première partie donne les trois questions à poser à n'importe quel
prélèvement. La deuxième les pose, un impôt après l'autre : l'argument en une phrase, un exemple
chiffré qu'on peut refaire de tête, les données internationales, la situation française, ce
qu'il faudrait faire et l'objection la plus forte. La troisième remonte au système d'ensemble.\par}
\vspace{12pt}
{\small\color{encre}Trois textes servent de fil : l'article de Mankiw, Weinzierl et Yagan (2009),
la \emph{Mirrlees Review} (2011) et le document de l'OCDE sur fiscalité et croissance (2008). Ils
sont cités là où ils éclairent un impôt particulier. Dix rapports publics français --- du Conseil
d'analyse économique, du Conseil des prélèvements obligatoires et de la Cour des comptes ---
confrontent ces raisonnements aux impôts français. Les données internationales viennent de l'OCDE
et de la Tax Foundation.\par}
\vspace{18pt}
{\color{bleu}\hrule height 1.6pt}

\clearpage
\thispagestyle{empty}
\vspace*{1.2cm}

{\small
{\fontsize{15}{18}\selectfont\bfseries\color{bleu}Sommaire\par}\vspace{14pt}
\begin{tabular}{@{}r@{\hspace{10pt}}l@{}}
& Avant de commencer : ce que la théorie peut dire, et les trois textes\\[6pt]
\multicolumn{2}{@{}l}{\itshape\color{ocre} Première partie --- Trois questions à poser à n'importe quel impôt}\\[2pt]
1 & Que détruit-il ?\\
2 & Qui le paie vraiment ?\\
3 & Comment le juger ?\\[6pt]
\multicolumn{2}{@{}l}{\itshape\color{ocre} Deuxième partie --- Impôt par impôt}\\[2pt]
4 & Les impôts sur la production\\
5 & L'impôt sur les sociétés\\
6 & Les prélèvements sur le travail\\
7 & La TVA\\
8 & L'épargne et le capital des ménages\\
9 & L'immobilier et les transmissions\\
10 & Le carbone\\[6pt]
\multicolumn{2}{@{}l}{\itshape\color{ocre} Troisième partie --- Le système d'ensemble}\\[2pt]
11 & Le niveau des prélèvements et le classement de l'OCDE\\
12 & Ce que font les autres pays\\
13 & Ce qu'il faudrait changer\\
\end{tabular}}

\clearpage

% ===========================================================================
\section*{Avant de commencer}
\addcontentsline{toc}{section}{Avant de commencer}

\subsection{Ce que la théorie peut dire}

La théorie de la taxation optimale ne dit pas combien il faut redistribuer. Le poids qu'on accorde
à un euro de plus pour un ménage modeste plutôt que pour un ménage aisé est un jugement de valeur,
comme le niveau souhaitable de la dépense publique. Les économistes qui prétendent le contraire font
passer leurs préférences pour des résultats.

Ce qu'elle établit tient en une phrase : pour un degré de redistribution donné et une recette
donnée, certains systèmes détruisent beaucoup moins de richesse que d'autres. Deux lecteurs qui ne
souhaitent pas la même progressivité peuvent donc s'accorder sur le diagnostic. Ni l'un ni l'autre
n'a intérêt à payer trente milliards de pertes pour obtenir ce que dix milliards suffiraient à
obtenir.

Si l'État pouvait observer ce que chacun est capable de gagner, la fiscalité ne poserait aucun
problème économique. Il suffirait de demander à chacun une somme fixe, calculée sur sa capacité et
non sur ce qu'il fait. Une somme qui ne dépend d'aucune décision ne décourage ni le travail, ni
l'épargne, ni l'investissement : on pourrait redistribuer autant qu'on le souhaite sans rien
détruire.

Le Royaume-Uni a essayé la version la plus simple de cette idée : la même somme pour tous, sans
rien chercher à savoir de la capacité de chacun. En 1989 en Écosse, l'année suivante en Angleterre et
au pays de Galles, il a remplacé l'impôt local assis sur la valeur locative des logements par une contribution
due par chaque adulte, la \emph{community charge}, vite appelée \emph{poll tax}. Elle n'avait
pourtant de forfaitaire que le principe. Chaque collectivité en fixait le montant, qui variait
fortement de l'une à l'autre. Les ménages modestes pouvaient en obtenir une réduction allant jusqu'à
80~\%, qui diminuait de 15 pence pour chaque livre de revenu supplémentaire : c'était un impôt de
plus sur le revenu, pour des ménages dont les aides baissaient déjà à mesure que leur revenu
augmentait.
L'assiette, enfin, était un registre auquel on pouvait se soustraire. Selon les estimations de
l'époque, jusqu'à 600\,000 personnes ont quitté les listes électorales pour ne pas être retrouvées,
et le non-paiement a dépassé la moitié des sommes dues dans plusieurs collectivités (Besley, Preston
et Ridge, 1997). La \emph{poll tax} a été remplacée dès 1993. L'expérience montre pourquoi un impôt
forfaitaire, efficient dans le modèle, ne se lève pas : une somme identique pour tous paraît
injuste dès qu'elle est élevée ; il faut alors la moduler selon le revenu, ce qui recrée la
distorsion qu'on voulait éviter ; et ceux qui ne peuvent pas payer cessent de payer.

L'obstacle de fond est ailleurs : l'État n'observe pas la capacité, il observe le revenu. Or le
revenu résulte à la fois de la capacité et de l'effort. Taxer le revenu pour atteindre ceux qui
peuvent beaucoup gagner revient à taxer aussi l'effort, donc à le décourager. Mirrlees (1971) l'a
formalisé et l'écrivait en toutes lettres : «~On pourrait obtenir des informations sur le
potentiel de gain d'un homme à partir de son quotient intellectuel apparent, du nombre de ses
diplômes, de son adresse, de son âge ou de sa couleur ; mais l'indicateur naturel, et
vraisemblablement le plus fiable, de son potentiel de gain est son revenu.~» Le reste du cours
découle de ce problème d'information.

\subsection{Les sources}

L'article de Mankiw, Weinzierl et Yagan, «~Optimal Taxation in Theory and Practice~»
(\emph{Journal of Economic Perspectives} 23(4), 2009), est le meilleur point d'entrée : vingt-huit
pages sans équation, huit leçons tirées de la théorie, chacune confrontée à l'évolution observée
des systèmes de l'OCDE. Exposer chaque résultat une fois comme conclusion d'un modèle, une fois
comme fait, évite la confusion la plus répandue dans les débats fiscaux.

La \emph{Mirrlees Review} (\emph{Tax by Design}, Institute for Fiscal Studies et Oxford
University Press, 2011, en accès libre) juge le système dans son ensemble plutôt qu'impôt par
impôt. Elle part d'une description chiffrée de la population à laquelle il s'applique et chiffre
chaque réforme avec ses perdants.

Le document de l'OCDE de Johansson, Heady, Arnold, Brys et Vartia, «~Taxation and Economic
Growth~» (document de travail n°~620, 2008), est le plus cité des trois. Il classe les impôts
selon leur effet sur la croissance à partir de régressions sur vingt et un pays ; ses auteurs
écrivent eux-mêmes que l'ampleur des effets estimés dépasse ce qu'on peut raisonnablement attendre.
Le chapitre~11 dit ce que vaut ce classement.

Dix rapports publics français appliquent ces raisonnements aux impôts français. Leur poids tient à
leurs données. Les notes du Conseil d'analyse économique, placé auprès du Premier ministre, sont
écrites par des économistes universitaires et n'engagent que leurs auteurs ; plusieurs contiennent
une estimation originale sur données d'entreprises ou données fiscales. Le Conseil des
prélèvements obligatoires, présidé par le Premier président de la Cour des comptes, s'appuie sur
des études qu'il commande à l'Insee, à l'Institut des politiques publiques ou à l'administration.
La Cour des comptes évalue des dispositifs à partir des fichiers de l'administration fiscale et
sociale. La liste figure à la fin du cours.

\partie{Première partie}{Trois questions à poser à n'importe quel impôt}

% ===========================================================================
\section{Que détruit-il ?}

Prenons une personne qui ferait volontiers deux heures de ménage par semaine chez un voisin pour
20~\euro{} de l'heure, et ce voisin, prêt à payer jusqu'à 26~\euro{} pour s'en dispenser. Sans
impôt, ils s'entendent par exemple sur 23~\euro{} : elle gagne 3~\euro{} de plus que le minimum
qu'elle exigeait, il paie 3~\euro{} de moins que le maximum qu'il acceptait. Chaque heure crée
6~\euro{} de valeur, du seul fait qu'ils se sont mis d'accord.

Un prélèvement de 10~\euro{} par heure fait disparaître cet échange. Pour qu'elle touche encore
20~\euro{}, il devrait débourser 30~\euro{}, et il refuse. L'État ne perçoit rien, puisqu'il n'y a
plus rien à taxer, et les 6~\euro{} de valeur ont disparu avec l'échange : personne ne les a.

Ces 6~\euro{} sont ce que les économistes appellent la perte sèche. Sur les échanges qui ont encore
lieu, l'impôt fait passer de l'argent des parties à l'État sans rien détruire. Sur ceux qu'il rend
non rentables, il détruit de la valeur sans rien rapporter. La première partie apparaît dans les
comptes publics ; la seconde consiste en ce qui ne s'est pas produit, et n'apparaît nulle part.

\subsection{Le coût croît avec le carré du taux}

Imaginons mille échanges possibles, tous différents : le plus avantageux rapporte 100~\euro{} à
ses deux parties, le moins avantageux presque rien, et les autres s'étagent régulièrement entre
les deux. Une taxe de 10~\euro{} par échange fait disparaître ceux qui rapportaient moins de
10~\euro{}. Il y en a cent, et ils rapportaient en moyenne 5~\euro{} : la perte vaut 500~\euro{}.
Une taxe de 20~\euro{} fait disparaître deux cents échanges, qui rapportaient en moyenne
10~\euro{} : la perte vaut 2\,000~\euro{}. Doubler la taxe a doublé le nombre d'échanges supprimés
et doublé la valeur moyenne de chacun ; la perte est multipliée par quatre. Les recettes, elles,
passent seulement de 9\,000 à 16\,000~\euro{}. Chaque euro prélevé coûte 5,6 centimes de perte
avec une taxe de 10~\euro{}, 12,5 avec une taxe de 20~\euro{}, 21,4 avec une taxe de 30~\euro{}
(figure~\ref{fig:carre}).

\figbloc{c1_perte_au_carre}{Doubler une taxe quadruple la perte}{fig:carre}

La règle vient de Jules Dupuit, ingénieur des Ponts et Chaussées, qui l'énonce en 1844 : «~l'utilité
perdue est proportionnelle au quarré de la taxe ; ainsi la taxe 10 fr. ferait perdre 100 fois plus
d'utilité que la taxe 1 fr.~» La règle découle du calcul : elle est exacte quand les échanges
menacés s'étagent régulièrement, comme dans l'exemple, et approchée sinon, tant que la taxe reste
modérée ; Dupuit la réservait lui-même aux «~taxes faibles~». On ne la teste donc pas directement. Ce que les études mesurent est la sensibilité des comportements à l'impôt,
dont on déduit ensuite le coût. Avec une élasticité du revenu imposable de 0,25, au milieu des
estimations disponibles (voir plus bas), la formule de Saez, Slemrod et Giertz (2012) donne le coût
d'un euro de recette supplémentaire : 7 centimes de perte s'il est obtenu en relevant un taux de
20~\%, 20 centimes à partir d'un taux de 40~\%, 60 centimes à partir de 60~\%. Le coût du dernier
euro prélevé croît plus vite que le taux : c'est la traduction pratique de la règle du carré.

Trois conséquences pratiques en découlent. À recette égale, une assiette large taxée à 20~\%
détruit deux fois moins qu'une assiette deux fois plus étroite taxée à 40~\%. Deux taux différents
sur des biens semblables coûtent plus qu'un taux unique qui rapporte autant : 10~\% sur l'un et
30~\% sur l'autre détruisent l'équivalent de $10^2 + 30^2 = 1\,000$, un taux de 20~\% sur les deux
$20^2 + 20^2 = 800$. Une niche fiscale, enfin, n'est jamais gratuite, même compensée ailleurs : elle
rétrécit l'assiette, oblige à relever le taux sur ce qui reste, et la perte croît plus vite que le
taux.

\subsection{Ce que la loi fait bouger}

Un impôt ne détruit de la valeur que dans la mesure où les gens peuvent y échapper en changeant de
comportement. Dans l'exemple du ménage, tout tenait au fait que le voisin pouvait renoncer. Si le
service lui avait été indispensable, il aurait payé 30~\euro{}, l'État aurait encaissé 10~\euro{}
et rien n'aurait été détruit. La sensibilité de l'assiette à l'impôt s'appelle son élasticité.

On la mesure par rapport à ce qui reste au contribuable sur un euro supplémentaire, et non par rapport
au taux d'impôt. Avec un taux marginal de 50~\%, il garde 50 centimes ; si le taux descend à 45~\%,
il en garde 55, soit 10~\% de plus. Une élasticité de 1 signifie que son revenu déclaré augmente
alors de 10~\%, une élasticité de 0,25 qu'il augmente de 2,5~\%. L'élasticité relie deux
pourcentages, celui du revenu et celui de ce qui reste après impôt ; un point d'impôt en plus ne
fait donc pas perdre un point de revenu.

Le revenu déclaré peut réagir de deux façons. Le contribuable peut travailler moins, investir moins,
renoncer à une promotion : c'est la réponse réelle, et aucune réforme de l'assiette ne la fait
disparaître. Il peut aussi déclarer autrement le même revenu, en le faisant passer dans une
catégorie moins taxée, en le reportant à une année plus favorable, en le logeant dans une société
ou un placement exonéré. Cette seconde réponse dépend de ce que la loi permet. En haut de la
distribution, c'est surtout elle qu'on observe : selon la synthèse de Saez, Slemrod et Giertz
(2012), les réactions mesurées sont pour l'essentiel des déplacements dans le temps et des montages
plutôt que des changements d'activité, et les meilleures estimations de l'élasticité vont de 0,12
à 0,40. En France, Cabannes, Houdré et Landais (2014) trouvent 0,02 pour l'ensemble des foyers et
0,31 pour les 10~\% les plus aisés, sur les réformes de 1997 à 2004 ; Lehmann, Marical et Rioux
(2013) trouvent environ 0,2 pour les salariés payés moins de deux fois le salaire minimum.

C'est ce qui donne à l'assiette son importance. Une loi qui ferme les possibilités de requalifier un
revenu supprime la seconde réponse et laisse la première intacte : l'élasticité qui reste est plus
faible, et le coût d'un taux donné aussi. Slemrod et Kopczuk (2002) en concluent que l'élasticité
est en partie un choix de politique publique.

\subsection{Où se trouve le sommet des recettes}

Quand un taux monte, chaque euro d'assiette rapporte davantage, mais il y a moins d'euros
d'assiette. Au-delà d'un certain taux, le second effet l'emporte et les recettes baissent. C'est la
courbe de Laffer, prise au sérieux : si l'élasticité est constante et vaut $e$, le sommet se
trouve exactement au taux $1/(1+e)$, soit 80~\% pour une élasticité de 0,25, 67~\% pour 0,5 et
50~\% pour 1 (figure~\ref{fig:sommet}). L'usage politique habituel de la courbe --- une baisse
d'impôt qui se financerait d'elle-même --- suppose qu'on se trouve au-delà de ce sommet. Pour un
taux de 30~\%, il faudrait une élasticité supérieure à 2,3, six à vingt fois les meilleures
estimations.

\figbloc{c12_sommet_recettes}{Le sommet des recettes ne dépend que de l'élasticité de
l'assiette}{fig:sommet}

Ce sommet est une borne, pas un objectif. Au-delà, monter le taux appauvrit à la fois les
contribuables et l'État ; aucun objectif ne justifie de s'y placer. En deçà, chaque point
supplémentaire rapporte, de moins en moins, et coûte de plus en plus à ceux qui le paient. Le taux
souhaitable se trouve donc entre zéro et le sommet, et sa position dépend du poids accordé au
revenu de ceux qui le paient. Si ce poids est nul, le taux souhaitable est au sommet. S'il est le
même que celui d'un euro de recette publique, le taux souhaitable est nul : prendre un euro à
quelqu'un pour le donner à l'État n'apporte alors rien, et la perte sèche s'y ajoute. Le choix de
ce poids est politique ; la position du sommet, elle, se calcule. Près de ce sommet, la courbe est
d'ailleurs plate : pour une élasticité de 0,5, les recettes restent à moins de 3~\% de leur maximum
entre 60 et 75~\% de taux.

% ===========================================================================
\section{Qui le paie vraiment ?}

Celui qui verse l'impôt au Trésor n'est pas forcément celui qui le supporte. La charge glisse vers
celui qui peut le moins s'y soustraire. Si l'on instaure une cotisation patronale de dix points,
l'employeur peut accepter de payer plus, modérer les salaires futurs ou embaucher moins. S'il y a
beaucoup de candidats pour peu de postes, c'est le salaire qui s'ajuste, et le salarié supporte la
cotisation que l'employeur verse. La loi désigne celui qui paie ; celui qui supporte dépend de qui
a le plus d'alternatives.

Les droits de mutation en donnent l'exemple le plus simple. Supposons l'offre de logements fixe à
court terme. Ce que l'acheteur est prêt à décaisser au total, prix et droits compris, ne change pas
parce qu'on instaure un droit : le prix doit donc baisser d'autant pour que ce total reste le même.
L'acheteur fait le virement à l'administration, le vendeur supporte l'impôt.

Une entreprise, elle, ne supporte jamais un impôt, parce qu'elle n'est qu'un ensemble de contrats
et n'a pas de bien-être propre. Tout prélèvement qu'elle acquitte se retrouve dans une baisse du
revenu de ses actionnaires, de ses salariés ou de ses fournisseurs, ou dans une hausse du prix
payé par ses clients. La répartition entre eux est souvent difficile à établir, y compris pour les
économistes, ce qui rend ces impôts politiquement commodes : ils paraissent sans victime. C'est
probablement aussi pourquoi la part patronale des cotisations sociales dépasse la part salariale
dans seize des dix-neuf pays de notre échantillon (dix-huit pays européens et les États-Unis) : un
prélèvement qui semble payé par l'employeur se vote plus facilement.

La comparaison entre la France et le Danemark le montre. L'OCDE mesure ce qu'elle appelle le coin
fiscal : la part du coût total d'un salarié qui part en cotisations et en impôt sur le revenu, pour
un célibataire payé au salaire moyen. En France, sur 100~\euro{} que ce salarié coûte à son
employeur, 47~\euro{} vont aux cotisations et à l'impôt, dont 27~\euro{} de cotisations patronales ;
il en garde 53. Au Danemark, il en part 36, dont 35 par l'impôt sur le revenu
(figure~\ref{fig:coin}). Il s'agit d'un taux moyen, au salaire moyen : la part prélevée monte avec
le salaire, jusqu'à 54~\% en France à 1,67 fois le salaire moyen (chapitre~6). La différence de
composition entre les deux pays tient largement à une convention : ce que la France prélève sous le
nom de cotisation patronale, le Danemark le prélève sous le nom d'impôt sur le revenu. Pour le
salarié, ce qui compte est l'écart entre ce que son travail coûte et ce qu'il touche, quelle que
soit l'étiquette.

\figbloc{s2_coin_fiscal}{Le coin fiscal sur le travail et sa composition, 2025}{fig:coin}

% ===========================================================================
\section{Comment le juger ?}

La progressivité se juge sur l'ensemble du système, pas impôt par impôt. Le débat public fait
l'inverse : une hausse de TVA y est dite injuste parce que la TVA pèse davantage, rapportée au
revenu, sur les ménages modestes, même lorsqu'elle finance une prestation qui leur rapporte plus
qu'elle ne leur coûte. Or ce qui compte pour un ménage est ce qu'il paie au total et ce qu'il reçoit
au total. Le Conseil des prélèvements obligatoires en donne un exemple chiffré dans son rapport de
février 2023 sur la TVA. Pour un même coût de 9,4~Md\euro{}, supprimer la TVA sur l'alimentation
relèverait le niveau de vie des 10~\% de ménages les plus modestes de 1,5~\% ; verser un chèque de
1\,620~\euro{} aux ménages éligibles au chèque énergie le relèverait de 7,3~\%. La baisse de TVA,
présentée comme la mesure sociale, l'est près de cinq fois moins que la prestation, et rien n'empêche de
financer la prestation par un impôt qui, pris isolément, n'est pas progressif. Juger chaque impôt
séparément conduit à écarter des réformes qui, compensation comprise, améliorent la situation des
plus modestes. C'est la règle la plus difficile à faire admettre, et celle dont dépendent plusieurs
propositions de la troisième partie. Une taxe sur le tabac, elle, est franchement
régressive ; elle ne vise pas la redistribution et n'a pas à la viser.

\subsection{Sur une vie plutôt que sur une année}

Le revenu d'une année mesure mal la position d'une personne. Il varie d'abord avec l'âge : faible
pendant les études et les premières années de travail, il monte jusqu'au milieu de la carrière,
puis baisse à la retraite. Une photographie prise une année donnée mêle des étudiants, des actifs au
sommet de leur carrière et des retraités, et une partie de ce qu'elle présente comme des écarts entre
riches et pauvres sépare en réalité des âges. Il varie ensuite d'une année sur l'autre, avec le
chômage, une prime, les résultats d'un indépendant.

Suivre les mêmes personnes pendant longtemps permet de mesurer ce second effet. L'Insee a classé
chaque année, à partir de leurs déclarations de revenus, les personnes âgées de 25 à 49 ans en 2003,
selon leur revenu avant impôt, retraites et allocations chômage comprises (figure~\ref{fig:mobilite}).
Seize ans plus tard, 63~\% de celles qui faisaient partie des 20~\% les plus aisés en 2003 y sont
encore, et 3,5~\% seulement sont passées parmi les 20~\% les plus modestes ; en bas, 62~\% sont
restées et 2~\% ont rejoint le cinquième le plus aisé. L'échelle des revenus bouge donc lentement en
France. Elle bouge tout de même : au bout de trois ans, une personne sur cinq a déjà quitté le
cinquième le plus aisé. Le tableau que cite la \emph{Mirrlees Review} pour le Royaume-Uni est bien
plus mobile --- 38~\% des plus riches de 1991 encore au sommet en 2008, 10~\% parmi les plus
pauvres ---, mais il porte sur des personnes de tous âges, enfants et retraités compris, et sur un
revenu hebdomadaire déclaré à une enquête. Une partie des passages du haut vers le bas tient
probablement au départ à la retraite, que ni la \emph{Review} ni la statistique britannique dont
elle tire le tableau ne détaillent. Les deux séries ne se comparent pas.

\figbloc{c14_mobilite}{Seize ans plus tard, deux tiers des plus aisés le sont encore}{fig:mobilite}

Même modérée, cette mobilité a deux conséquences concrètes. La première est qu'un barème progressif
appliqué année par année taxe davantage un revenu irrégulier qu'un revenu stable de même total.
Prenons deux célibataires qui gagnent chacun 500\,000~\euro{} en dix ans. Le premier a un revenu
imposable de 50\,000~\euro{} chaque année et paie, au barème de 2025, 8\,165~\euro{} par an. Le
second gagne alternativement 20\,000 et 80\,000~\euro{} : il paie 470~\euro{} les mauvaises années
et 17\,165~\euro{} les bonnes, soit 8\,818~\euro{} en moyenne, 8~\% de plus pour le même revenu. Les
revenus irréguliers sont ceux des indépendants, des agriculteurs, des artistes, des salariés qui
traversent des périodes de chômage, et le barème annuel les pénalise sans que personne l'ait
décidé. Le droit français ne lisse que les revenus exceptionnels, par le système du quotient, qui
impose un tel revenu comme s'il avait été perçu sur quatre ans ; les variations ordinaires ne sont
pas lissées.

La seconde conséquence porte sur le diagnostic. Mesurées sur le revenu moyen de 2003 à 2020 plutôt
qu'année par année, les inégalités entre les personnes suivies par l'Insee sont inférieures de
6,6~\%. L'écart est modeste pour des actifs d'âge comparable ; il grandit nécessairement quand la
photographie mêle des âges différents, puisque s'y ajoutent les écarts entre jeunes, actifs et
retraités. C'est le cas de toutes les statistiques annuelles sur les ménages, et c'est ce qui fausse
le jugement porté sur la TVA.

Rapportée au revenu d'une année, la TVA paraît régressive : elle représente 12,5~\% du revenu disponible des 10~\% de ménages les plus modestes et
4,7~\% de celui des 10~\% les plus aisés (CPO, 2023). Mais une partie de ces ménages modestes ne le
sont que pour un temps --- étudiants, personnes entre deux emplois, retraités --- et consomment plus
qu'ils ne gagnent, en empruntant ou en puisant dans leur épargne. La figure~\ref{fig:tvavie} suit
une même personne à trois âges. Jeune, elle dépense plus qu'elle ne gagne ; au milieu de sa
carrière, elle épargne ; retraitée, elle dépense son épargne. La TVA représente pour elle 18~\%,
puis 12~\%, puis 20~\% du revenu de l'année, et une photographie qui mêlerait des personnes des
trois âges ferait apparaître un impôt régressif. Sur sa vie entière, elle paie 15~\% de ce qu'elle
gagne : la TVA est proportionnelle à ses ressources. D'où la règle de mesure que recommande la
\emph{Mirrlees Review} : rapporter un impôt sur la dépense à la dépense, et non au revenu de
l'année. Sur données françaises, une analyse sur le cycle de vie réalisée pour le CPO réduit la
régressivité de la TVA de près de moitié sans l'annuler, notamment parce que les ménages aisés ne
consomment pas toutes leurs ressources de leur vivant (chapitre~7).

\figbloc{c22_tva_cycle_de_vie}{La TVA paraît régressive sur une année, elle est proportionnelle
sur une vie}{fig:tvavie}

\subsection{Juger une réforme}

Une réforme ne s'évalue qu'avec sa compensation. Proposer un taux unique de TVA sans dire comment
le choc est compensé pour les bas revenus n'est que la moitié d'une proposition, et c'est ce qui
permet à une bonne idée d'être rejetée pour de mauvaises raisons. La \emph{Mirrlees Review} a
construit ses réformes pour qu'elles soient neutres à la fois sur la recette et sur la
distribution ; c'est ce qui leur a donné leur crédit.

La complexité, de son côté, profite à ceux qui peuvent payer un conseil. Les niches, les régimes
dérogatoires et les montages servent d'abord ceux qui savent les trouver ; à barème identique, un
système simple redistribue davantage qu'un système complexe.

La Cour des comptes chiffre ce que coûtent ces exceptions. Une dépense fiscale est une exception à
la règle générale d'un impôt, qui fait le travail d'une subvention sans passer par le budget ; la
Cour en recense 474 en 2025, pour 91,8~Md\euro{}, plus du quart des recettes fiscales nettes de
l'État (\emph{Le budget de l'État en 2025}, avril 2026). Elles représentent l'équivalent de 46~\%
du produit de l'impôt sur le revenu et de 52~\% de celui de l'accise sur les énergies. Pour 13~\%
d'entre elles le coût est inconnu, pour 42~\% le nombre de bénéficiaires. Le total dépend en outre
d'une convention : depuis 2024, les dépenses fiscales de TVA ne sont comptées que sur la part qui
revient à l'État, ce qui les divise par deux, et avec la convention antérieure le total de 2025
dépasse 105~Md\euro{}. La prévision, enfin, est systématiquement inférieure au coût définitif :
celui du pacte Dutreil a été multiplié par six pour 2024 après l'évaluation que la Cour en a faite
(chapitre~9). Le tableau~\ref{tab:niches} donne les plus coûteuses.

\begin{table}[htbp]
\caption{Les quinze dépenses fiscales les plus coûteuses, chiffrage pour 2026}
\label{tab:niches}
\footnotesize
\renewcommand{\arraystretch}{1.18}
\begin{tabularx}{\linewidth}{@{}r>{\raggedright\arraybackslash}X l r@{}}
\toprule
 & \textbf{Dépense fiscale} & \textbf{Impôt} & \textbf{Md\euro{}}\\
\midrule
1 & Crédit d'impôt recherche & IR et IS & 8,0\\
2 & Crédit d'impôt pour l'emploi d'un salarié à domicile & IR & 7,2\\
3 & Abattement de 10~\% sur les pensions et retraites & IR & 4,7\\
4 & Pacte Dutreil : transmission d'entreprises & successions & 4,0\\
5 & Participation, intéressement et épargne salariale & IR & 2,9\\
6 & TVA à 10~\% sur les travaux dans les logements & TVA & 2,5\\
7 & TVA à 10~\% sur la restauration & TVA & 2,3\\
8 & Exonération des heures supplémentaires & IR & 2,3\\
9 & Réduction d'impôt pour les dons des particuliers & IR & 2,2\\
10 & Déduction des travaux sur les logements loués & IR & 2,0\\
11 & Crédit d'impôt pour la garde des enfants de moins de 6 ans & IR & 1,8\\
12 & Réduction d'impôt pour les dons des entreprises & IR et IS & 1,7\\
13 & Gazole à tarif réduit pour les travaux agricoles et forestiers & accise & 1,7\\
14 & Exonération des prestations familiales et de l'allocation aux adultes handicapés & IR & 1,6\\
15 & Accise réduite sur les carburants outre-mer & accise & 1,5\\
\midrule
 & \textbf{Total des quinze} & & \textbf{46,2}\\
\bottomrule
\end{tabularx}

\vspace{4pt}
\parbox{\linewidth}{\raggedright\scriptsize\color{encre}Source : projet de loi de finances pour
2026, \emph{Évaluation des voies et moyens}, tome II, p.~25. Ces quinze mesures font plus de la
moitié du coût total des 474 dépenses fiscales. Le chiffrage du pacte Dutreil a été relevé de
0,8 à 4~Md\euro{} pour 2025 après l'évaluation de la Cour des comptes.}
\end{table}

La Cour relève aussi que douze dépenses fiscales dont le gain moyen par bénéficiaire est inférieur
à 100~\euro{} par an coûtent ensemble 2,2~Md\euro{}. Son rapport n'en donne pas la liste. En
appliquant le même critère aux chiffrages de 2023 publiés avec le projet de loi de finances pour
2025, on retrouve exactement douze dépenses pour 2,2~Md\euro{}, ce qui laisse penser qu'il s'agit
des mêmes. Les plus coûteuses sont l'exonération des revenus de l'épargne salariale (590~M\euro{},
soit 49~\euro{} en moyenne pour 12~millions de bénéficiaires), celle de la contribution de
l'employeur aux titres-restaurant (479~M\euro{}, 92~\euro{}), l'abattement accordé aux personnes
âgées ou invalides de condition modeste (362~M\euro{}, 52~\euro{}), la réduction d'impôt pour frais
de scolarité dans le secondaire (218~M\euro{}, 70~\euro{}) et l'exonération de la prise en charge
des frais de transport domicile-travail (189~M\euro{}, 33~\euro{}). Ces montants ne comptent que
l'impôt sur le revenu : certaines de ces exonérations valent aussi pour les cotisations sociales,
et l'avantage réel est alors plus élevé.

\subsection{Ce qu'on demande à un système fiscal}

Tout le monde s'accorde sur les quatre règles d'Adam Smith : que chacun paie selon ses facultés,
que l'impôt soit certain et non arbitraire, commode à payer et peu coûteux à percevoir. Elles ne
disent rien de ce qu'il faut faire quand elles se contredisent --- quand l'impôt le plus juste est
aussi le plus coûteux à percevoir, par exemple. La \emph{Mirrlees Review} pose une question plus
étroite : pour une redistribution donnée, quel système réduit au minimum les pertes d'activité, les
coûts de gestion pour l'administration et de déclaration pour les contribuables, tout en restant
compréhensible et prévisible ? La simplicité, la neutralité et la stabilité y servent ces
objectifs : un système simple coûte moins à gérer et se comprend mieux. La même logique autorise
une exception quand elle se justifie, comme le prix du carbone ou les accises sur le tabac, et
interdit d'invoquer la neutralité comme un argument qui se suffirait à lui-même.

\begin{encadre}{Le test du statu quo}
Ne pas taxer quelque chose revient à en subventionner l'usage par rapport à tout ce qui est taxé.
Un taux réduit sur la restauration est une subvention aux repas pris à l'extérieur, financée par
un taux plus élevé sur tout le reste. L'absence d'impôt sur le loyer implicite du propriétaire
occupant --- le loyer qu'il ne paie pas parce qu'il est chez lui --- est une subvention à la
propriété par rapport à la location.

La question n'est donc jamais «~faut-il taxer ceci ?~» mais «~pourquoi subventionne-t-on ceci ?~».
Posée dans ce sens, elle appelle une justification là où l'habitude n'en demande aucune, et la
plupart des régimes dérogatoires n'en ont pas.
\end{encadre}

On peut supposer, sans que la littérature l'ait démontré, que la subvention au propriétaire
occupant freine la croissance de long terme : l'épargne qu'elle attire vers le logement occupé ne
finance pas le capital des entreprises. Les modèles calibrés vont dans ce sens, et trouvent que
taxer le loyer implicite augmenterait le bien-être (Gervais, 2002) ; aucune étude ne mesure de façon
convaincante l'effet dans les données. Seuls quatre pays de l'OCDE taxent ce loyer implicite, et la
Suisse a voté en 2025 la suppression du sien. Sur le marché du travail,
l'hypothèse d'Oswald veut qu'une part élevée de propriétaires augmente le chômage, parce qu'un
propriétaire déménage moins facilement pour un emploi. La preuve la plus solide vient de Finlande,
où la libéralisation des loyers au début des années 1990 permet d'isoler l'effet : une hausse de la
part des propriétaires y a fait monter le chômage, alors même que les propriétaires sont eux-mêmes
moins souvent au chômage que les locataires (Laamanen, 2017). L'effet passerait par le marché local
dans son ensemble plutôt que par les propriétaires (Blanchflower et Oswald, 2013). La subvention à
la propriété fausse donc le choix entre placements, et probablement le fonctionnement du marché du
travail ; on ne sait pas en chiffrer le coût en croissance, et l'argument pour la supprimer reste
un argument de neutralité.

\begin{encadre}{Imposer l'individu ou le couple}
La France impose le foyer. Le revenu d'un couple marié ou pacsé est divisé en deux parts ; l'impôt
est calculé sur une part, puis multiplié par deux. Chaque enfant ajoute une demi-part, le troisième
une part entière. Trois exemples, au barème de 2025, montrent ce que ce choix change.

Un couple dont un seul membre gagne 60\,000~\euro{} de revenu imposable paie 4\,331~\euro{}.
Imposés comme deux personnes seules, ils paieraient 11\,165~\euro{} : le quotient conjugal leur fait
gagner 6\,835~\euro{}, et cet avantage n'est pas plafonné. Un couple dont chaque membre gagne
30\,000~\euro{} paie 4\,331~\euro{} dans les deux systèmes : le quotient ne profite qu'aux couples
dont les revenus sont inégaux. Si, dans le premier couple, le second conjoint prend un emploi qui
lui rapporte 15\,000~\euro{}, l'impôt du foyer augmente de 4\,500~\euro{} : son salaire est taxé à
30~\% dès le premier euro, alors qu'imposé seul il ne paierait rien.

Le système favorise donc le mariage et le PACS des couples dont les revenus sont inégaux --- les
couples non mariés sont imposés séparément --- et décourage le travail du second apporteur de
revenu, une femme dans les trois quarts des cas. Selon l'Insee,
l'imposition commune relève de 5,9 points le taux marginal des seconds apporteurs et abaisse de
13 points celui des premiers. L'imposition individuelle ne favorise ni ne pénalise le mariage, et
taxe chaque salaire à son propre taux. On peut supposer, sans que la littérature l'ait établi pour
la France, que l'avantage fiscal du mariage et du PACS encourage la formation de couples stables,
ce qu'un État peut juger souhaitable ; il bénéficie toutefois surtout aux couples où l'un des deux
gagne peu ou pas. Pour les enfants, la demi-part réduit l'impôt de
1\,791~\euro{} par an au plus, et d'autant plus que le revenu est élevé ; son effet mesuré sur la
natalité est positif mais très faible (Landais, 2003). Un crédit d'impôt d'un même montant par
enfant donnerait autant pour chaque enfant ; le quotient repose sur un autre principe, qui compare
les familles à niveau de vie égal.

Aucun résultat théorique ne départage l'imposition de l'individu et celle du couple, et la
\emph{Mirrlees Review} renonce à trancher. Le Conseil des prélèvements obligatoires chiffre le coût
du quotient conjugal entre 10,6 et 13,5~Md\euro{}. L'avantage peut atteindre 35\,355~\euro{} pour un
foyer, et près de la moitié du total va aux 10~\% de ménages les plus aisés. À l'inverse, trois
millions de couples, pour l'essentiel des classes moyennes, paient plus que s'ils étaient imposés
séparément, parce que la décote qui allège l'impôt des petits contribuables n'est pas doublée pour
un couple. Le Conseil propose de la doubler, pour un coût de 1,3~Md\euro{}, et de plafonner
l'avantage du quotient conjugal à 10\,000~\euro{}, ce qui rapporterait 1,1~Md\euro{}.
\end{encadre}

\partie{Deuxième partie}{Impôt par impôt}

La France prélève 43,9~\% du PIB en 2023, dix points de plus que la moyenne des pays de l'OCDE
(33,7~\%). L'écart se concentre sur quelques catégories : 5,7 points viennent des cotisations
sociales et 1,5 des taxes sur la masse salariale ; suivent l'impôt sur le revenu, l'impôt foncier et
les successions (figure~\ref{fig:ecarts}). Une seule grande catégorie rapporte moins que la
moyenne : l'impôt sur les sociétés, 2,35~\% du PIB contre 3,75~\%. Il s'agit de recettes, pas de
taux. Le taux français est un peu au-dessus de la moyenne (chapitre~5), mais ce qu'il rapporte
dépend aussi de la part des bénéfices dans l'économie et de l'étroitesse de l'assiette, et la
moyenne est tirée vers le haut par des pays qui taxent des rentes de ressources naturelles
(Norvège, 12~\% du PIB ; Colombie, Australie, Chili) ou qui accueillent les bénéfices de
multinationales (Luxembourg, Pays-Bas). L'Allemagne et les États-Unis, à 2,2~\%, sont au niveau de
la France ; la médiane de l'OCDE est à 3,4~\%.

\figbloc{r1_carte_des_ecarts}{D'où viennent les dix points d'écart avec la moyenne de
l'OCDE}{fig:ecarts}

% ===========================================================================
\section{Les impôts sur la production}

\begin{argument}
Un impôt sur ce qui sert à produire réduit la production avant qu'on ait commencé à la partager.
On peut toujours lever la même somme, sur les mêmes personnes, sans cette perte : c'est le seul
gain d'efficience qui ne coûte rien à personne.
\end{argument}

Un intrant est ce qu'une entreprise utilise pour produire : du travail, des machines, des locaux,
des fournitures achetées à d'autres entreprises. Taxer un intrant, comme les locaux ou la masse
salariale, ou taxer les ventes que les entreprises se font entre elles, change la façon dont
l'entreprise combine ses moyens. Elle utilise moins
de l'élément taxé et davantage des autres, uniquement parce que c'est fiscalement moins coûteux, et
l'économie produit moins avec les mêmes ressources. C'est le théorème d'efficience productive de
Diamond et Mirrlees (1971) : l'État peut taxer ce que les ménages consomment, ce qu'ils gagnent, ce
qu'ils possèdent, mais il ne devrait pas taxer ce qui s'échange entre entreprises pour produire.

Ce résultat a une propriété rare en fiscalité : il ne demande aucun arbitrage entre efficience et
équité. Si l'on supprime une taxe sur un intrant et qu'on lève le même montant sur la consommation
ou le revenu des mêmes personnes, chacun se retrouve dans la situation d'avant, avec davantage à
partager.

Le théorème suppose toutefois que l'État peut taxer en aval tout ce qu'il veut taxer : la
consommation, les salaires, les dividendes. Quand il ne le peut pas, prélever au niveau de
l'entreprise redevient utile, et c'est la meilleure justification de l'impôt sur les sociétés. Un
entrepreneur qui laisse ses bénéfices dans sa société, ou dans une société holding, ne se verse ni
salaire ni dividende : son revenu échappe à l'impôt sur le revenu aussi longtemps qu'il le
souhaite, et seul l'impôt sur les sociétés le frappe l'année où il est gagné. L'impôt sur les
sociétés atteint aussi les actionnaires étrangers, hors de portée de l'impôt sur le revenu
français (chapitre~5). Il n'est d'ailleurs pas un impôt sur un intrant au sens du théorème : il
frappe le bénéfice, ce qui reste une fois tous les intrants payés. Les impôts de ce chapitre, eux,
frappent les intrants avant tout bénéfice.

\subsection{L'exemple : une taxe de 0,16 \% qui en vaut 0,40}

La contribution sociale de solidarité des sociétés (C3S) prélève 0,16~\% du chiffre d'affaires des
sociétés, au-delà d'un abattement de 19~millions d'euros de ventes ; elle a rapporté 5,3~Md\euro{}
en 2025. Un prélèvement assis sur les ventes frappe le même produit chaque fois qu'il change de mains.
Supposons un produit vendu 100~\euro{} au consommateur et fabriqué en quatre étapes qui ajoutent
chacune 25~\euro{} de valeur. Une entreprise qui ferait tout elle-même paierait la taxe une fois,
sur 100~\euro{}. Dans une chaîne de quatre entreprises, la première la paie sur 25~\euro{}, la
deuxième sur 50, la troisième sur 75, la quatrième sur 100 : la taxe porte au total sur
250~\euro{}, et elle représente 0,40~\% du prix final au lieu de 0,16
(figure~\ref{fig:pyramidage}). Faire fabriquer une pièce par un sous-traitant coûte donc plus cher
que la fabriquer soi-même, sans que personne l'ait décidé. La TVA ne produit pas cet effet, puisque
chaque entreprise déduit la taxe déjà payée par son fournisseur : c'est précisément pour supprimer
ces cascades qu'elle a été inventée.

\figbloc{c7_pyramidage}{Une taxe sur le chiffre d'affaires frappe plusieurs fois le même
produit}{fig:pyramidage}

L'Allemagne en a fait l'expérience. Jusqu'en 1967, elle prélevait une taxe de 4~\% sur le chiffre
d'affaires à chaque stade, dont son ministre des Finances écrivait en 1963 qu'elle «~favorise la
concentration verticale~». En 1966, la Cour constitutionnelle fédérale a constaté qu'elle n'était
pas neutre entre entreprises intégrées et non intégrées, tout en l'admettant jusqu'à la réforme alors
en préparation : la TVA l'a remplacée le 1\ier{}~janvier 1968. La France avait connu le même défaut
avec sa taxe sur le chiffre d'affaires de 1920, dont la charge dépendait de la longueur du circuit de
production ; la TVA de Maurice Lauré, en 1954, est l'aboutissement des réformes engagées pour le
supprimer. L'effet a pu être mesuré là où l'expérience a été faite proprement : dans l'État de
Washington, remplacer une taxe de 25~\% prélevée à chaque stade de la filière du cannabis par une
taxe sur la seule vente au détail a réduit l'intégration verticale et augmenté la production de
23~\% (Hansen, Miller et Weber, 2022). À 0,16~\%, la C3S est bien plus faible, et aucune étude ne
mesure son effet sur l'organisation des filières françaises ; le sens de l'effet, lui, n'est pas
douteux. On peut supposer, sans que la littérature l'ait montré, qu'elle freine aussi les effets
d'agglomération : les entreprises proches les unes des autres gagnent à se spécialiser et à
s'acheter mutuellement ce que chacune fait le mieux, et c'est précisément cet échange que la taxe
renchérit à chaque étape. Selon le CAE, aucun autre pays
européen n'a d'équivalent ; les plus proches sont des taxes sur le chiffre d'affaires de quelques
États américains.

Personne ne se mobilise contre un taux de 0,16~\%. Mais ce taux s'applique au chiffre d'affaires,
et il est dû que l'entreprise gagne ou perde de l'argent. Pour un distributeur dont le bénéfice
avant impôt représente 1~\% des ventes, 0,16~\% du chiffre d'affaires en prend 16~\% ; à 0,5~\% de
marge, près d'un tiers ; en perte, la taxe reste due. Le législateur l'a admis pour quelques métiers
seulement : dans le commerce international, le négoce de produits agricoles et celui des
combustibles et carburants, la contribution est plafonnée à 3,08~\% de la marge brute lorsque celle-ci ne dépasse
pas 4~\% du chiffre d'affaires, et pour les banques à 1,6~\% de leur produit net bancaire. Les
autres activités à faible marge n'ont pas de plafond.

La taxe pousse aussi à s'organiser pour l'éviter. L'abattement de 19~millions s'applique société
par société, sans consolidation au niveau du groupe et sans règle particulière contre le
fractionnement : un groupe qui répartit ses ventes entre plusieurs filiales plus petites paie moins.
Un intermédiaire qui vend en qualité de commissionnaire, sans devenir propriétaire des
marchandises, n'est taxé que sur sa commission ; un distributeur qui achète et revend l'est sur tout
le prix. Quand ils sont faits pour la taxe, ces choix d'organisation ne produisent rien et coûtent du
conseil.

La note du Conseil d'analyse économique de juin 2019, signée par Philippe Martin et Alain Trannoy,
a mesuré ce pyramidage sur l'économie française. Elle suit la taxe le long des échanges entre
branches décrits par la comptabilité nationale, en supposant que chaque entreprise la répercute à
ses clients. Dans l'industrie manufacturière, la taxe relève les prix d'environ 0,19~\%, près de
deux fois son taux effectif de 0,11~\% ; selon les secteurs, le rapport va de 1,2 à 3. Les auteurs
comparent ensuite près de 80\,000 entreprises manufacturières : celles que les abattements de 2015
et 2016 ont fait sortir de l'assiette ont vu leurs exportations progresser d'environ 1~\% de plus
que les autres, résultat qu'ils qualifient eux-mêmes de préliminaire. La taxe n'est pas remboursée
à l'exportation, et un bien intermédiaire importé n'a pas traversé les étapes de la cascade qu'a
traversées un bien produit en France : elle agit, selon le mot des auteurs, comme un droit de douane
que la France imposerait à sa propre production.

Un modèle multisectoriel chiffre enfin la seule perte de productivité à 360 à 720~millions d'euros
de PIB par an, 10 à 20~\% de ce que la taxe rapportait alors. Le montant paraît modeste ; c'est le
rapport à la recette qui frappe. Une taxe de 0,16~\% sur un bien final ne détruirait, d'après la
règle du carré du chapitre~1, qu'une fraction de centime par euro prélevé. Le modèle trouve 10 à 20
centimes, parce que la taxe porte sur des intrants et se propage d'entreprise en entreprise. Ce
chiffre ne compte ni l'effet sur les exportations ni les défaillances, que la même note estime plus
fréquentes dans les entreprises assujetties pendant la crise de 2009.

\subsection{Ce que montrent les données}

Sur les 1\,275~Md\euro{} de prélèvements obligatoires français en 2024, 119,3~Md\euro{} sont
calculés sur un intrant, c'est-à-dire sur quelque chose que l'entreprise utilise pour produire ---
sa masse salariale, ses locaux, son chiffre d'affaires --- plutôt que sur ce qu'elle vend aux
ménages ou sur son bénéfice. Parmi eux, 42,2 corrigent une externalité : ce sont pour l'essentiel
les taxes sur l'énergie, que le théorème admet comme exception. Les 77,0 restants ne corrigent rien.
La taxe sur les salaires en représente 17,3, le versement mobilité 12,2, les contributions à la
formation professionnelle et à l'apprentissage 11,4, la cotisation foncière des entreprises 7,7, la
contribution de solidarité 5,2 (figure~\ref{fig:77md}).

\figbloc{r11_77_milliards}{Les prélèvements français assis sur un facteur de production et sans
fonction corrective}{fig:77md}

Ces 77 milliards ne forment pas un bloc. La contribution de solidarité et la cotisation foncière
frappent la production elle-même. Les taxes assises sur la masse salariale --- taxe sur les
salaires, versement mobilité, formation, logement, autonomie --- sont des cotisations sans droits :
elles n'ouvrent aucune prestation proportionnelle aux sommes versées, et finissent largement dans
des salaires plus bas (chapitre~2). Les premières devraient disparaître ; les secondes devraient
être fondues dans le prélèvement général sur les revenus du travail, plutôt que maintenues comme
autant d'impôts sur l'emploi à assiette étroite. Dans les deux cas, la recommandation est de lever
les mêmes sommes sur des assiettes larges.

Sur les taxes assises sur la masse salariale, la France est troisième de l'OCDE : elles atteignent
1,9~\% du PIB, quatre fois la moyenne de 0,46~\%, derrière la Suède et l'Autriche
(figure~\ref{fig:masse}).

\figbloc{r5_masse_salariale}{Les taxes sur la masse salariale, hors cotisations sociales, dans
l'OCDE}{fig:masse}

Les annonces de baisse n'ont pas fait reculer l'ensemble en trente ans. La catégorie que la
comptabilité européenne appelle «~autres impôts sur la production~» valait 4,1 points de PIB en
1995 ; elle en vaut 4,4 en 2024, malgré la suppression progressive de la cotisation sur la valeur
ajoutée engagée en 2021 (figure~\ref{fig:production}). Chaque suppression a été compensée par une
création ailleurs, ou par la croissance des bases restantes. Le même agrégat vaut 1,0 point en
Allemagne et 1,2 aux Pays-Bas, mais son contenu varie d'un pays à l'autre. La note du CAE contourne
la difficulté en retirant les taxes assises sur la masse salariale : ce qui reste représentait en
France 3,6~\% de la valeur ajoutée des entreprises, contre 0,5~\% en Allemagne ; l'Autriche, les
Pays-Bas et la Suède prélevaient entre 0,5 et 1,5~\%, et seule la Grèce davantage que la France.

\figbloc{c17_production_depuis_1995}{Les impôts sur la production n'ont pas reculé en trente
ans}{fig:production}

\subsection{Ce qu'il faudrait faire}

La taxe sur le chiffre d'affaires vient en premier, parce qu'elle cumule les deux défauts : elle
frappe un intrant et elle se multiplie le long de la chaîne. Deux rapports publics y arrivent à six
ans d'écart. Le CAE recommandait en 2019 de finaliser la suppression de la C3S, engagée en 2015 puis
interrompue en 2017, avant de programmer celle de la cotisation sur la valeur ajoutée des
entreprises (CVAE). Le Conseil des prélèvements obligatoires, dans son rapport de septembre 2025 sur
l'industrie, juge que si l'effort en faveur des entreprises se poursuit, il doit porter sur la C3S
plutôt que sur l'achèvement de la suppression de la CVAE, prévu pour 2030 : l'industrie, qui
produit 14,5~\% de la valeur ajoutée, acquitte 29,6~\% de la C3S. Il propose de la financer en
supprimant les exonérations fiscales et sociales des heures supplémentaires, qui coûtent un montant
comparable et n'ont pas d'effet mesuré sur les heures travaillées (chapitre~6).

La CVAE vient ensuite. Son assiette, la valeur ajoutée, est neutre en principe : c'est la somme de
ce que l'entreprise verse à ses salariés et de ce qui lui reste, et la taxer ne pousse pas à
préférer les machines aux personnes. Mais son taux dépendait du chiffre d'affaires, si bien qu'à
valeur ajoutée égale le CAE trouvait des taux allant du simple au double, et elle ne déduit pas les
achats de machines, qu'elle taxe donc à leur tour. Faut-il alors la remplacer par un impôt sur le
bénéfice ? En partie. Un impôt sur le bénéfice n'est dû que si l'entreprise gagne de l'argent, et il
ne se multiplie pas le long de la chaîne de production ; le CAE le juge moins nocif que les impôts
sur la production. Il a cependant ses propres défauts : il frappe le rendement normal du capital, ce
qui renchérit l'investissement, et il porte sur des bénéfices que les groupes peuvent déplacer d'un
pays à l'autre (chapitre~5). Le CAE proposait donc d'en élargir
l'assiette, par une meilleure imposition des multinationales et la réduction de ses niches, plutôt
que d'en relever le taux, et de supprimer des taux réduits de TVA peu efficaces, comme celui de la
restauration. La suppression se finance d'ailleurs en partie d'elle-même : la CVAE étant déduite du
bénéfice imposable, sa disparition augmente mécaniquement l'impôt sur les sociétés, de
2,6~Md\euro{} selon le chiffrage de 2019. La recette des collectivités serait remplacée par une
part d'un impôt national.

La cotisation foncière des entreprises taxe en partie le capital productif, mais le CAE n'y voit pas
de distorsion majeure, et la part de sa base qui correspond au terrain pourrait être reprise par un
impôt foncier assis sur le sol (chapitre~9). Les taxes sur la masse salariale viennent en dernier,
pour la raison donnée plus haut.

La taxe sur les salaires demande un traitement à part. Elle est due par les employeurs qui ne
collectent pas la TVA --- banques, assurances, établissements de santé, associations --- et
fonctionne donc comme un substitut grossier à la TVA que ces secteurs n'acquittent pas. La supprimer
sans rien mettre à la place leur donnerait un avantage fiscal. La \emph{Mirrlees Review} recommande de soumettre les services
financiers à un équivalent de la TVA. La façon la plus simple, proposée par le FMI en 2010, consiste
à taxer la somme des salaires et des bénéfices, c'est-à-dire la valeur ajoutée calculée par
addition. Prenons une banque dont les clients sont des ménages : ses marges d'intérêts et ses
commissions lui rapportent 120, elle achète 30 à d'autres entreprises, verse 50 de salaires et
dégage 40 de bénéfice. Sa valeur ajoutée vaut 90 ; une taxe de 20~\% sur les salaires et le bénéfice
rapporte 18, exactement ce que rapporterait une TVA de 20~\% sur ses marges, 24, diminuée de la TVA
payée sur ses achats, 6. Aujourd'hui, exonérée de TVA, elle ne verse que ces 6, qu'elle ne peut pas
récupérer, et une taxe sur les salaires qui ignore le bénéfice et renchérit l'emploi plutôt que le
capital. Israël applique ce principe aux banques, en taxant leurs salaires et leurs bénéfices au
taux de la TVA ; le Danemark ne taxe que la masse salariale du secteur financier. La difficulté
tient aux clients professionnels : une entreprise ne pourrait pas déduire la taxe contenue dans ses
frais bancaires, et la cascade reviendrait. Il faut donc retirer de l'assiette les opérations
réalisées avec des entreprises, ce qui complique le calcul ; la \emph{Review} ne choisit pas entre
les méthodes possibles, mais recommande d'en adopter une.

\subsection{L'objection la plus forte}

Ces prélèvements financent des politiques précises --- transports urbains, formation, logement,
autonomie --- et beaucoup sont affectés à un organisme ou à une collectivité. La règle de
l'universalité budgétaire, qui interdit d'affecter une recette à une dépense, ne vaut que pour le
budget de l'État : le versement mobilité va aux autorités organisatrices des transports, les
contributions à la formation à France compétences, la cotisation au Fonds national d'aide au
logement à ce fonds. Supprimer la ressource oblige à en trouver une autre pour chacun de ces
organismes, et chacun a de bonnes raisons de défendre la sienne.

L'objection vise pourtant la dépense, que la réforme ne touche pas. Rien n'oblige à financer les
transports urbains par un impôt sur l'emploi plutôt que par un impôt sur le foncier qu'ils
valorisent, ni la formation par une taxe sur les salaires plutôt que par une dotation. L'argent
étant fongible, transférer le financement vers une assiette large à recette égale laisse chaque
politique intacte et supprime la perte. L'affectation rend la réforme plus difficile à négocier ;
elle ne la rend pas moins juste.

On peut aussi supposer, sans que la littérature l'ait prouvé, qu'un impôt dû même sans bénéfice
a un effet utile : pousser les entreprises à abandonner leurs activités les moins rentables et à
monter en gamme. Les travaux disponibles ne le confirment pas, et le mécanisme s'y prête mal. Selon la note
du CAE, 20~\% des entreprises redevables de
la C3S ne font pas de bénéfice, et la contribution a accru la probabilité de disparition des
entreprises assujetties pendant la crise de 2009, sans indication qu'elle ait frappé les moins
productives. Elle ne le peut pas : une taxe sur le chiffre d'affaires trie les entreprises selon
leur marge, pas selon leur productivité. Un grossiste dont la valeur ajoutée représente 6~\% des
ventes paie en C3S l'équivalent de 2,7~\% de sa valeur ajoutée ; un industriel dont la valeur
ajoutée représente 35~\% des ventes, 0,46~\%. Le premier n'est pas moins productif, il exerce un
métier où l'on vend beaucoup pour une petite marge, et le législateur l'a reconnu en plafonnant la
contribution de certains négociants. L'argument a une version sérieuse, mais pour un autre impôt :
un impôt sur le patrimoine, dû quel que soit le rendement, pèse relativement plus sur les
entrepreneurs dont le capital rapporte peu (Guvenen \emph{et al.}, 2023).

% ===========================================================================
\section{L'impôt sur les sociétés}

\begin{argument}
L'impôt sur les sociétés est le seul moyen d'atteindre les bénéfices que leurs propriétaires
laissent dans l'entreprise et ceux qui reviennent à des actionnaires étrangers. Tel qu'il est
construit, il frappe aussi le rendement normal de l'investissement et favorise l'endettement. Son
assiette compte davantage que son taux.
\end{argument}

\subsection{L'exemple : cent euros de bénéfice}

Une société qui réalise 100~\euro{} de bénéfice paie d'abord l'impôt sur les sociétés, au taux
normal de 25~\% : il lui reste 75~\euro{}. Si elle les distribue à un actionnaire qui réside en
France, celui-ci acquitte le prélèvement forfaitaire unique de 30~\% --- 12,8~\% d'impôt sur le
revenu et 17,2~\% de prélèvements sociaux ---, soit 22,5~\euro{}. Il garde 52,5~\euro{} : le
bénéfice a supporté 47,5~\% d'impôt en deux étapes (figure~\ref{fig:cascade}). En 2026, la hausse
de la CSG sur les revenus de placement porte le prélèvement forfaitaire à 31,4~\%, et ce taux
combiné à 48,6~\%. Une petite entreprise, dont le chiffre d'affaires ne dépasse pas 10~millions
d'euros et dont le capital est détenu aux trois quarts par des personnes physiques, ne paie que
15~\% sur ses 42\,500 premiers euros de bénéfice.

Trois contributions s'ajoutent pour les grandes entreprises et les gros actionnaires. Une
entreprise dont l'impôt dépasse 763\,000~\euro{}, soit un bénéfice d'environ 3~millions d'euros,
paie une contribution sociale de 3,3~\% de l'impôt au-delà de ce montant, ce qui porte son taux
marginal à 25,8~\%. Un actionnaire dont le revenu fiscal de référence dépasse 250\,000~\euro{}, pour
une personne seule, paie la contribution exceptionnelle sur les hauts revenus, de 3 puis 4~\% ; ce
revenu de référence comprend les dividendes. Le taux combiné atteint alors 51~\% ; parmi les pays de
la figure, seuls l'Irlande, le Danemark et le Royaume-Uni font davantage. Depuis 2025, enfin, la
contribution différentielle sur les hauts revenus oblige les foyers dont le revenu de référence
dépasse ce seuil à payer au moins 20~\% de ce revenu en impôt sur le revenu et en contribution sur
les hauts revenus. Pour un foyer qui vit surtout de dividendes, elle porte le prélèvement sur le
dividende à 37,2~\% et le taux combiné à 53,4~\%.

Le taux de 36,1~\% souvent cité pour la France ne concerne que les plus grandes entreprises. La loi
de finances pour 2025 a créé une contribution exceptionnelle sur leurs bénéfices, égale à 20,6~\% de
l'impôt pour celles dont le chiffre d'affaires réalisé en France dépasse un milliard d'euros et à
41,2~\% au-delà de trois milliards ; celle de 2026 l'a reconduite pour un an aux mêmes taux, en
relevant le premier seuil à 1,5~milliard. Le taux d'impôt passe ainsi à 31,0~\% dans le premier cas
et à 36,1~\% dans le second, et le taux combiné au sommet à 57,8~\%, le plus élevé des pays
comparés. Le Gouvernement prévoyait environ 450 entreprises redevables en 2025, et la contribution a
rapporté 7,5~Md\euro{} cette année-là. Pour toutes les autres entreprises, le taux reste de 25 ou
25,8~\%.

La double imposition décrit qui verse l'impôt. Qui en supporte la charge dépend d'autre chose. La
société verse l'impôt sur les sociétés, mais celui-ci réduit ce que rapporte, après impôt, un investissement
réalisé en France. Si les investisseurs peuvent obtenir ailleurs le même rendement, ils n'investissent
en France que dans les projets assez rentables pour compenser l'impôt. Il y a alors moins de
machines par salarié, la production par salarié est plus faible, et les salaires aussi : une partie
de l'impôt finit chez des salariés qui n'en ont jamais versé un euro. Ce n'est pas qu'une hypothèse. En Allemagne, où chaque commune fixe le taux de
son impôt sur les bénéfices des entreprises, Fuest, Peichl et Siegloch (2018) ont suivi 6\,800
changements de taux sur vingt ans : les salariés supportent environ la moitié de la charge. Aux
États-Unis, Suárez Serrato et Zidar (2016) l'attribuent pour 40~\% aux actionnaires, pour 30 à
35~\% aux salariés et pour 25 à 30~\% aux propriétaires du sol. Le taux combiné
garde pourtant son sens. Il mesure l'écart entre ce que l'investissement rapporte avant impôt et ce
qu'en reçoit l'actionnaire, et c'est cet écart qui décide si l'investissement se fait ; la façon
dont la charge se répartit ensuite entre actionnaires et salariés est la question de l'incidence,
traitée au chapitre~2.

\figbloc{c6_cascade_profit}{Ce que deux impôts successifs font à cent euros de bénéfice}{fig:cascade}

Les deux impôts se réforment donc ensemble. Baisser l'impôt sur les sociétés en relevant d'autant
le prélèvement sur les dividendes ne change rien pour l'actionnaire résident qui reçoit tout le
bénéfice en dividende. Cela allège en revanche la charge sur les bénéfices réinvestis dans
l'entreprise, et sur ceux qui reviennent aux actionnaires étrangers. Ces derniers ne paient pas le
prélèvement français sur les dividendes : ils sont imposés dans leur pays, la France ne prélevant
qu'une retenue à la source que les conventions fiscales limitent en général. Pour la part du
bénéfice qui leur revient, l'impôt sur les sociétés est le seul impôt français. C'est l'un de ses
arguments les plus solides, et c'est aussi pourquoi la concurrence fiscale entre pays porte sur
son taux : c'est lui que voit l'investisseur étranger.

\subsection{Le mécanisme : rendement normal et rente}

Pour calculer son bénéfice imposable, une entreprise déduit ses coûts. Une machine qui servira vingt
ans n'est pas déduite l'année de son achat : son coût est étalé, ici 5~\euro{} par an pendant vingt
ans pour une machine de 100~\euro{}. Or un euro déduit dans dix ans vaut moins qu'un euro déduit
aujourd'hui, comme un euro reçu dans dix ans : à 5~\% d'intérêt, il vaut 61 centimes. Additionnées
en valeur d'aujourd'hui, les vingt déductions de 5~\euro{} ne valent que 62~\euro{}
(figure~\ref{fig:amort}). L'entreprise est imposée comme si la machine ne lui avait coûté que
62~\euro{} ; 38~\euro{} de son coût ne sont jamais déduits. Plus l'inflation et les taux d'intérêt
sont élevés, plus l'écart se creuse, sans qu'aucun taux affiché ait bougé.

\figbloc{c8_amortissement}{Ce que vaut aujourd'hui une déduction étalée dans le
temps}{fig:amort}

Ce coût non déduit fait peser l'impôt sur ce que les économistes appellent le rendement normal :
ce que l'investissement doit rapporter simplement pour rémunérer l'argent immobilisé, au taux qu'on
obtiendrait ailleurs. La rente est ce qui dépasse --- un emplacement unique, un brevet, une position
de marché. Prenons un investissement de 100~\euro{} qui rapporte exactement 5~\euro{} par an quand
l'actionnaire pourrait obtenir 5~\% ailleurs. Taxé, il rapporte moins de 5~\% après impôt et ne
sera pas fait, alors qu'il crée de la valeur. Un investissement qui rapporte 15~\euro{} sera fait
même taxé, puisqu'il lui reste une partie des 10~\euro{} de rente. Un impôt qui ne frappe que la rente
ne décourage aucun investissement ; un impôt qui frappe aussi le rendement normal décourage ceux qui
sont tout juste rentables.

L'impôt sur les sociétés favorise en outre l'endettement. Les intérêts d'emprunt sont déductibles,
le coût des fonds propres ne l'est pas : une entreprise qui finance un investissement en empruntant
paie moins d'impôt que celle qui le finance avec l'argent de ses actionnaires. La \emph{Mirrlees
Review} propose de traiter les fonds propres comme la dette, en laissant l'entreprise déduire
chaque année un rendement normal sur ses capitaux propres : c'est la déduction pour fonds propres.
Dans l'exemple, une entreprise financée par 100~\euro{} de fonds propres déduirait 5~\euro{} par an.
Si l'investissement rapporte 5~\euro{}, son bénéfice imposable est nul ; s'il rapporte 15~\euro{},
elle est imposée sur 10~\euro{}, c'est-à-dire sur la rente seulement. Elle est traitée comme si
elle s'était endettée, et le biais en faveur de la dette disparaît.

La même déduction répare le défaut de l'amortissement. Tant que la machine n'est pas entièrement
amortie, la part de son coût qui n'a pas encore été déduite reste dans les fonds propres de
l'entreprise et ouvre droit, chaque année, à la déduction de 5~\%. Plus l'amortissement est lent,
plus cette part reste élevée longtemps, et plus la déduction pour fonds propres est importante. Les
deux effets se compensent exactement : dans l'exemple de la machine amortie sur vingt ans, les
déductions pour fonds propres valent 38~\euro{} aujourd'hui, ce qui manquait pour revenir à
100~\euro{}. Quel que soit le rythme d'amortissement, l'entreprise déduit donc l'équivalent de ce
qu'elle a investi, et l'impôt ne frappe plus que la rente.

\subsection{Ce que montrent les données}

Le taux français était de 50~\% en 1985, un peu au-dessus de la moyenne des pays de l'OCDE
(47~\%). Il est resté entre 34 et 38~\% de 2006 à 2019, pendant que la moyenne passait de 27 à
24~\%, et n'est descendu à 25,8~\% qu'en 2022 : la France a suivi la baisse générale avec une
dizaine d'années de retard (figure~\ref{fig:is_temps}). La contribution exceptionnelle de 2025
porte le taux des plus grandes entreprises à 36,1~\%, le plus élevé de l'OCDE, mais c'est le taux
de 25,8~\% que paient les autres.

\figbloc{r2_is_depuis_1985}{Le taux d'impôt sur les sociétés depuis 1985}{fig:is_temps}

La baisse du taux s'est retrouvée dans ce que les entreprises paient. Rapporté à l'excédent brut
d'exploitation, qui est l'assiette économique de l'impôt, l'impôt sur les sociétés des entreprises
industrielles est passé, avant crédits d'impôt, de 23,2 à 17,5~\% entre 2016 et 2022, et celui de
l'ensemble des entreprises de 20,7 à 17,5~\%, selon une étude de la DGFiP et de l'Insee réalisée
pour le Conseil des prélèvements obligatoires.

Un taux élevé ne fait pas pour autant une recette élevée. En 2023, le taux français se situait au
treizième rang sur trente-huit, sa recette, 2,35~\% du PIB contre 3,75~\% en moyenne, au septième
rang en partant du bas (figure~\ref{fig:is_recette}). Le graphique n'établit aucune relation
causale : la recette dépend aussi de la part des bénéfices dans l'économie, de la taille du secteur
des sociétés et des régimes particuliers.

\figbloc{r3_is_taux_et_recettes}{Taux statutaire et recettes de l'impôt sur les sociétés dans
l'OCDE}{fig:is_recette}

La même dissociation apparaît dans le temps. Dans les dix-huit pays de l'OCDE pour lesquels les
séries sont complètes depuis 1981, le taux légal moyen de l'impôt sur les sociétés a été divisé par
deux, de 48 à 24~\%, pendant que sa recette passait de 1,9 à 3,5~\% du PIB
(figure~\ref{fig:tauxrecettes}). La France a suivi le mouvement des taux avec retard, et sa recette
a peu augmenté.

\figbloc{c23_taux_et_recettes_is}{Le taux de l'impôt sur les sociétés a été divisé par deux, sa
recette a presque doublé}{fig:tauxrecettes}

Des travaux plus fins mesurent la réaction de l'assiette au taux, en comparant des entreprises
soumises à des changements de taux différents. Une méta-analyse de Beer, de Mooij et Liu (2020)
conclut qu'une baisse d'un point du taux d'un pays augmente d'environ 1~\% les bénéfices que les
filiales des multinationales y déclarent, et de 1,5~\% dans les années récentes : une partie de la
base est déplacée d'un pays à l'autre plutôt que créée ou détruite. Sur données fiscales
américaines, Coles, Patel, Seegert et Smith (2022) trouvent qu'une hausse de 10~\% du taux marginal
anticipé réduit le revenu imposable des sociétés non cotées de 9,1~\%, dont 3 points de réaction
réelle et 6 points de réaction comptable. Tørsløv, Wier et Zucman (2023), par une méthode
comptable qui n'identifie pas de causalité, estiment que 36~\% des bénéfices des multinationales
étaient logés dans des paradis fiscaux en 2015 et que la France y perdait 21~\% de son impôt sur les
sociétés. Le taux qui maximiserait la recette reste mal connu : 33~\% selon Clausing (2007) sur
les pays de l'OCDE de 1979 à 2002, beaucoup plus dès qu'on tient compte des différences permanentes
entre pays (Kawano et Slemrod, 2016) ; ce sont des corrélations.

La valeur actuelle des déductions permet de comparer les règles d'amortissement des pays. La Tax
Foundation calcule,
pour 100~\euro{} investis, ce que valent aujourd'hui les déductions que chaque pays autorise, en
actualisant à 7,5~\% par an : 100 correspondrait à une déduction immédiate, et plus la déduction est
étalée, plus le chiffre est bas. Pour la France, elle retient un amortissement sur vingt ans pour
un bâtiment industriel : 5~\euro{} de déduction par an, dont la somme vaut 55~\euro{} aujourd'hui.
Un actif incorporel acquis, comme un brevet, est amorti sur cinq ans : 20~\euro{} par an, qui valent
87~\euro{} parce qu'ils arrivent plus tôt. Pour une machine, le chiffre est de 88~\euro{}. La France
est au-dessus de la moyenne de l'OCDE pour les trois actifs, qui est de 50 pour les bâtiments, 77
pour les incorporels et 86 pour les machines (figure~\ref{fig:amort_ocde}). Le défaut décrit plus
haut est réel, mais il est commun à presque tous les pays et ne distingue pas la France.

\figbloc{r4_amortissements}{La valeur actuelle des amortissements dans les pays de
l'OCDE}{fig:amort_ocde}

Ce qui la distingue est ce qu'elle prélève avant le bénéfice. Le taux effectif moyen sur un
investissement, qui tient compte du taux, de l'assiette et des amortissements, s'établit à 23,8~\%
en France, contre 27,0~\% en Allemagne, 25,4~\% en Italie et 24,2~\% en Belgique. Les impôts sur la
production du chapitre précédent, eux, sont acquittés que l'entreprise gagne de l'argent ou non :
ils représentent 4,4~\% du PIB en France, contre 1,0~\% en Allemagne (figure~\ref{fig:entreprises}).

\figbloc{s3_entreprises}{Imposition des entreprises : le bénéfice et ce qui le précède}{fig:entreprises}

\subsection{Ce qu'en disent les trois textes}

Mankiw, Weinzierl et Yagan relèvent que la baisse des taux affichés ne dit pas tout de la charge
sur le capital. Les taux légaux de l'impôt sur les sociétés sont passés de 45 à 50~\% au milieu des
années 1980 à environ 28~\% en moyenne dans l'OCDE en 2007, et le taux moyen de l'impôt personnel
sur les dividendes de 55~\% à moins de 20~\%. Carey et Rabesona rapportent pourtant l'ensemble des
impôts payés sur les revenus du capital à ces revenus, tels que les mesure la comptabilité
nationale, et ce ratio a augmenté : de 39,9~\% en 1975-1980 à 46,3~\% en 1990-2000 en moyenne dans
seize pays de l'OCDE, de 42,0 à 55,9~\% en France (figure~\ref{fig:carey}). La recette de l'impôt
sur les sociétés a suivi le même chemin jusqu'à aujourd'hui (figure~\ref{fig:tauxrecettes}).

\figbloc{c24_carey_rabesona}{La charge effective sur le capital a augmenté pendant que les taux
baissaient}{fig:carey}

Ces chiffres ne mesurent pas le seul effet de l'élargissement des assiettes. La recette dépend aussi
de la part des bénéfices dans le revenu national et de la part de l'activité exercée en société
plutôt qu'en entreprise individuelle. Le passage à la forme sociétaire, encouragé par l'écart croissant entre
l'impôt sur les sociétés et l'impôt sur le revenu, expliquerait à lui seul 12 à 21~\% des recettes
européennes de l'impôt sur les sociétés (de Mooij et Nicodème, 2008). L'élargissement des assiettes
est néanmoins réel. Dans plus de la moitié des années où un pays de l'OCDE a changé son taux entre
1980 et 2004, il a aussi changé son assiette, et les baisses de taux se sont plus souvent
accompagnées d'un élargissement que d'un rétrécissement (Kawano et Slemrod, 2016). Devereux,
Griffith et Klemm (2002) en montrent l'effet : le taux effectif sur l'investissement tout juste
rentable est resté à peu près stable, celui sur les investissements très rentables a baissé. On
peut donc supposer, sans que la littérature ait réussi à séparer les effets, qu'une part
importante de la hausse des recettes tient à la combinaison d'un taux plus bas et d'une assiette
plus large, qui fait payer davantage pour un taux affiché moindre. Les ratios de Carey et Rabesona sont en outre
sensibles aux conventions : si l'on attribue au travail une partie du revenu des indépendants, la
hausse moyenne est réduite d'environ deux points, et la hausse française de 14 points tombe à moins
de 2. Mankiw, Weinzierl et Yagan en tirent une conclusion prudente : la baisse des taux légaux ne
prouve pas que l'imposition du capital ait baissé. Les gouvernements ont pu appliquer la règle selon
laquelle un taux bas sur une assiette large décourage moins l'activité qu'un taux élevé sur une
assiette étroite, ou réduire la partie la plus visible de l'impôt en compensant ailleurs. Dans les
deux cas, un débat conduit sur les seuls taux affichés manque l'essentiel.

Le document de l'OCDE ajoute deux recommandations qui concernent directement la France. Pour
attirer l'investissement, baisser le taux de l'impôt sur les sociétés vaut mieux que baisser
l'imposition des dividendes, qui ne touche que les actionnaires résidents. Mais un taux d'impôt sur
les sociétés très inférieur au taux supérieur de l'impôt sur le revenu compromet l'intégrité du
système : les contribuables à haut revenu logent alors leurs revenus dans des sociétés qui ne
distribuent pas. En France, l'écart atteint trente points --- 25,8~\% contre 55,4~\% ---, et le
chapitre~8 montre ce qu'il produit dans les holdings patrimoniales.

\subsection{Ce qu'il faudrait faire}

Il faudrait un taux stable, sans surtaxe, et une assiette qui ne frappe que la rente. La déduction
d'un rendement normal sur les fonds propres supprime à la fois l'impôt sur le rendement normal et le
biais en faveur de la dette. Elle réduit l'assiette ; pour garder la même recette, il faut appliquer
un taux un peu plus élevé à ce qui reste, c'est-à-dire à la rente, qui ne se décourage pas. La Tax
Foundation compte en 2025 cinq pays de l'OCDE qui ne taxent pas le rendement normal des fonds
propres : la Pologne, le Portugal et la Turquie, par une déduction de ce type, l'Estonie et la
Lettonie, parce qu'elles n'imposent les bénéfices qu'au moment où ils sont distribués.

La stabilité est aussi ce que demande le Conseil des prélèvements obligatoires. Son rapport de
septembre 2025 propose d'inscrire dans la loi de programmation des finances publiques une
trajectoire des prélèvements par grandes catégories, la suppression de la contribution
exceptionnelle sur les bénéfices des grandes entreprises et le maintien du taux normal à 25~\%. Il
propose de financer la fin de la contribution en partie en resserrant les dépenses fiscales qui
réduisent le rendement de l'impôt, comme le crédit d'impôt recherche, qu'il voudrait calculer au
niveau du groupe, ou les régimes réservés à certaines zones. L'argument tient à la durée de
l'investissement industriel : une usine se décide pour quinze ou vingt ans, sur la base du taux qui
s'appliquera pendant toute cette durée, et une entreprise qui ne peut pas le prévoir attend, ou
exige un rendement plus élevé avant d'investir.

\subsection{L'objection la plus forte}

Une déduction pour fonds propres coûte cher, et les deux pays européens qui l'avaient appliquée le
plus largement y ont renoncé. La
Belgique l'a introduite en 2006 sur la totalité des fonds propres, au taux des obligations d'État à
dix ans, entre 3 et 4,5~\% selon les années. Les entreprises belges ont renforcé leurs fonds propres :
leur part dans le bilan des entreprises moyennes et grandes a augmenté d'environ trois points par
rapport à des entreprises comparables des pays voisins (Meki, 2023). Mais le coût a explosé. Annoncé à
un demi-milliard d'euros par an, il a atteint 5,4~Md\euro{} en 2011 et 6,2~Md\euro{} en 2012 selon
l'administration fiscale belge, plus de la moitié de la recette de l'impôt sur les sociétés, en
partie parce que des groupes ont logé en Belgique des fonds propres qu'ils prêtaient ensuite à leurs
filiales étrangères : environ 10~Md\euro{} pour les seuls groupes allemands (Hebous et Ruf, 2017).
La Belgique a plafonné le taux, l'a limitée à la hausse des fonds propres en 2018, puis l'a supprimée
en 2023. L'Italie avait choisi d'emblée cette version incrémentale, qui ne porte que sur les fonds
propres apportés depuis 2010. Elle a réduit d'environ 9 points l'endettement des entreprises
bénéficiaires (Branzoli et Caiumi, 2020), pour un coût resté la première année sous les prévisions,
avant d'être supprimée en 2024 pour financer d'autres mesures, contre l'avis de l'office
parlementaire du budget. Sur l'investissement, les résultats sont faibles et contradictoires : une
hausse de 3~\% pour les PME belges et de 6~\% pour les filiales belges de multinationales selon deux
études, aucun effet selon une troisième.

L'objection porte sur la mise en œuvre plus que sur le principe. Une déduction sur le stock de
fonds propres coûte cher dès le premier jour, parce qu'elle récompense aussi les fonds propres qui
existaient déjà et ne changent aucune décision ; une déduction sur les seuls fonds propres nouveaux
coûte peu au départ, et son coût monte à mesure qu'elle touche plus d'investissements. Dans les deux
cas, il faut des règles contre les montages entre filiales. Le FMI chiffre le coût d'une déduction
complète à environ 15~\% de la recette de l'impôt sur les sociétés, 0,5~\% du PIB, dont les
trois quarts seraient récupérés à long terme selon ses simulations. Celles de
de Mooij et Devereux (2011) trouvent un gain pour l'économie si la déduction est financée par des
impôts sur le travail ou la consommation, et une perte si elle l'est par une hausse du taux de
l'impôt sur les sociétés, qui fait fuir les bénéfices. La réforme ne vaut donc que financée par une
assiette moins mobile. Quant à l'impôt minimum mondial de 15~\%, il fixe un plancher pour les
grands groupes, pas un taux souhaitable, et ne dit rien de l'imposition du rendement normal d'une
usine qui peut être construite ailleurs.

% ===========================================================================
\section{Les prélèvements sur le travail}

\begin{argument}
Ce qui décourage le travail n'est pas ce qu'on paie en moyenne, mais ce qu'on paie sur l'euro
suivant. En France, cet euro est le plus lourdement prélevé au bas de l'échelle, pas au sommet,
et le niveau moyen du prélèvement est élevé dès le salaire minimum.
\end{argument}

\subsection{L'exemple : taux moyen et taux marginal}

Le taux moyen rapporte ce que l'on paie à ce que l'on gagne. Le taux marginal est ce que l'on paiera sur l'euro suivant. Le premier détermine ce qu'il vous reste. Le second
détermine, seul, vos décisions : une heure de plus, une promotion, un revenu déclaré sous une
forme plutôt qu'une autre. Un salarié dont le taux marginal atteint 60~\% ne perd pas 60~\% de
son revenu ; son taux moyen peut rester à 20~\%. Toute la théorie du barème optimal porte sur le
taux marginal, presque toute la discussion publique sur le taux moyen.

La progressivité se définit sur le taux moyen. Elle n'exige donc pas que le taux marginal augmente. Un impôt est
progressif lorsque le taux moyen croît avec le revenu. Un taux unique de 30~\% au-delà d'un
abattement de 12\,000~\euro{} l'est donc : celui qui gagne 20\,000~\euro{} a un taux moyen de
12~\%, celui qui gagne 100\,000~\euro{} de 26,4~\%, et tous deux ont le même taux marginal
(figure~\ref{fig:flat}).

\figbloc{c11_flat_tax_progressive}{Un taux unique assorti d'un abattement est déjà un impôt
progressif}{fig:flat}

Trois leviers rendent un tel impôt plus progressif, à des coûts différents.
Relever l'abattement ne change le taux marginal de personne au-dessus du seuil : c'est le plus
efficient. Relever le taux unique le change pour tout le monde. Ajouter une tranche supérieure ne
le change que pour ceux qu'elle atteint, mais c'est précisément là que les revenus imposables
réagissent le plus.

\subsection{Ce que dit la théorie : un barème en U}

Le taux marginal optimal à un niveau de revenu donné dépend de trois choses. L'élasticité du revenu imposable : plus elle est forte, plus le taux doit être bas. Le
nombre de personnes situées au-dessus de ce niveau, rapporté au nombre de personnes situées
exactement à ce niveau. Et le poids que l'on accorde au revenu de ceux qui sont au-dessus, seul
paramètre réellement politique.

Le deuxième ingrédient explique toute la forme du barème. Un taux marginal appliqué à un
niveau de revenu rapporte sur \emph{tous} ceux qui gagnent davantage : pour eux, cet euro-là est
déjà franchi, et l'impôt supplémentaire est acquis sans qu'ils changent quoi que ce soit. Il ne
coûte en distorsion que sur ceux qui sont \emph{exactement} à ce niveau, puisqu'eux seuls
arbitrent sur cet euro précis. En bas, presque tout le monde est au-dessus et très peu de gens sont
exactement là : un taux élevé est optimal --- c'est ce que fait mécaniquement tout transfert sous
condition de ressources, et c'est le prix du ciblage. Au milieu, où se trouve le gros de la
population, le taux doit baisser. En haut, il remonte, parce que ceux qui sont au-dessus
concentrent une masse importante de revenu.

Au sommet, le taux qui maximise les recettes dépend presque entièrement d'un paramètre mal connu. La formule vaut $1/(1+a\,e)$, où $e$ est l'élasticité du revenu
imposable et $a$ l'épaisseur de la queue haute de la distribution, autour de 2 en France. Avec une
élasticité de 0,25, elle donne 67~\% ; avec 0,5, elle donne 50~\% ; avec 1, elle donne 33~\%
(figure~\ref{fig:taupsup}). L'écart entre deux estimations plausibles de l'élasticité couvre
l'essentiel du désaccord politique sur le sujet.

\figbloc{c2_taux_superieur}{Le taux marginal supérieur qui maximise les recettes}{fig:taupsup}

L'histoire de ce paramètre montre comment une hypothèse de forme commande une conclusion
de politique publique. Mirrlees (1971) obtenait un taux marginal nul tout en haut : un taux
positif sur le dernier euro de la personne la plus riche décourage son effort sans rien
rapporter, puisque personne n'est au-dessus. Tuomala (1990), en supposant une distribution
log-normale des capacités, obtenait des taux \emph{décroissants} sur une large plage de hauts
revenus. Diamond (1998) et Saez (2001), en supposant une queue de Pareto, plus épaisse, obtenaient
l'inverse : des taux croissants, et au sommet des taux «~qui ne devraient pas être inférieurs à
50~\% et peuvent atteindre 80~\%~». Aucune donnée nouvelle ne sépare ces conclusions, seulement le
choix de la loi ajustée à la queue de la distribution --- et Mankiw, Weinzierl et Yagan montrent
que, sur les salaires américains correspondant à des revenus annuels de 100\,000 à
200\,000~dollars, les deux ajustements sont visuellement indiscernables.

\begin{encadre}{Trois avertissements sur le taux supérieur}
\textbf{Premièrement}, le taux calculé maximise les \emph{recettes} ; il n'est le taux
\emph{optimal} que si l'on accorde un poids nul au bien-être des contribuables concernés.

\textbf{Deuxièmement}, l'élasticité dépend de l'assiette que la loi a construite. Le barème au
sommet se discute avec l'assiette qui le porte, jamais séparément.

\textbf{Troisièmement}, si les gens diffèrent aussi par leur goût pour le travail ou la
consommation, et pas seulement par leur capacité, l'étendue optimale de la redistribution
diminue, parce qu'une part de l'inégalité observée reflète alors des choix.
\end{encadre}

Deux résultats de Mankiw, Weinzierl et Yagan tirent en sens opposés. Le premier est qu'un impôt à taux unique assorti d'un transfert universel peut
être proche de l'optimum : Mirrlees lui-même relevait «~la quasi-linéarité des barèmes~» obtenus.
Ce qui redistribue n'est donc pas le nombre de tranches, mais le niveau du transfert et celui du
taux. Le second est que l'ampleur optimale de la redistribution croît avec l'inégalité des
salaires : un sommet plus riche rend un euro prélevé en haut moins coûteux en bien-être et finance
davantage en bas. Sur les onze pays de la \emph{Luxembourg Income Study} pour lesquels la série est
exploitable, neuf montrent une relation positive entre inégalité avant impôt et dépenses
sociales --- une corrélation, pas une identification causale, mais la prédiction du modèle n'est
pas contredite.

\subsection{Travailler ou travailler plus}

En bas de l'échelle, la décision change de nature. Un
cadre arbitre sur l'intensité de son travail. Une personne au voisinage du minimum social arbitre
souvent sur la participation : prendre un emploi, ou pas. Quand la réaction porte sur cette
seconde marge, ce qui compte n'est plus le taux marginal mais le gain total que procure l'emploi.
Saez (2002) a montré que le barème optimal subventionne alors le premier euro de revenu
d'activité au lieu de le taxer : c'est la logique de l'\emph{Earned Income Tax Credit} américain
et, en France, de la prime d'activité.

Le document de l'OCDE tire de cette distinction une règle qui n'est jamais passée dans le
débat français. Réduire les prélèvements moyens sur le travail favorise la participation ;
réduire les taux marginaux favorise les heures travaillées. Ce ne sont pas les mêmes instruments.
La France a massivement utilisé le premier --- les allègements de cotisations au voisinage du
salaire minimum --- et ce faisant a dégradé le second, puisque la sortie progressive de ces
allègements crée précisément des taux marginaux élevés.

La Cour des comptes a chiffré cette pente. Son rapport
de mai 2025 sur la sécurité sociale montre que les cotisations patronales passent de 4,3~\% du
salaire brut au salaire minimum à 43,4~\% à 3,5 fois ce salaire. Les allègements qui creusent cet écart ont
coûté 77,3~Md\euro{} en 2024, presque quatre fois les 20,9~Md\euro{} de 2014. Les neuf
revalorisations du salaire minimum intervenues depuis 2021 expliquent une grande partie de la
hausse, par un effet que personne n'avait voté : les plafonds sont exprimés en multiples du salaire
minimum, si bien que chaque revalorisation fait entrer sous les plafonds des salaires qui n'ont pas
augmenté aussi vite.

Ce que ces allègements ont produit dépend du niveau de salaire où ils portent. La note
du CAE de janvier 2019, signée par Yannick L'Horty, Philippe Martin et Thierry Mayer, rassemble les
évaluations existantes et en ajoute une sur données d'entreprises. Les baisses concentrées sous
1,6 fois le salaire minimum ont un effet positif sur l'emploi, que les évaluations retrouvent
depuis les allègements des années 1990. Celles qui ont été étendues jusqu'à 3,5 fois ce salaire en
2016, au nom de la compétitivité, n'ont eu aucun effet détectable sur les exportations, alors que
ces salaires représentent 10~\% de la valeur de la production des exportateurs. Une explication
possible, selon les auteurs, relève du chapitre~2 : à ces niveaux de qualification, le chômage est
faible, les employeurs se disputent les mêmes salariés, et la baisse du coût finit dans leurs
salaires. Le Conseil des prélèvements obligatoires la retient en 2025 : au-delà du voisinage du
salaire minimum, les allègements augmentent le salaire net sans améliorer directement la
compétitivité. Le crédit d'impôt pour la compétitivité et l'emploi, enfin, a créé ou sauvegardé de
l'ordre de 100\,000 emplois pour un coût près de six fois supérieur à celui d'un allègement ciblé
au voisinage du salaire minimum dont les auteurs attendaient 80\,000 à 200\,000 emplois ; la
comparaison rapproche une évaluation après coup et une estimation avant coup, et ils la présentent
comme un ordre de grandeur.

L'exonération des heures supplémentaires montre ce que la loi fait bouger quand elle
abaisse un taux marginal. Elle réduit précisément le prélèvement sur l'heure en plus, et elle
coûte de l'ordre de 5~Md\euro{} en 2025. Les évaluations que cite le Conseil des prélèvements
obligatoires constatent une hausse des heures déclarées comme supplémentaires, sans effet
significatif sur les heures travaillées. La marge qui a réagi est la déclaration, pas le travail :
c'est l'élasticité fabriquée par la loi du chapitre~1.

\subsection{Ce que montrent les données}

Le taux supérieur français ne concerne presque personne. En 2025, il
atteint 55,4~\%, prélèvements assimilés compris, mais ne s'applique qu'à partir de treize fois le
salaire moyen. La Belgique (52,7~\%) et le Danemark (55,9~\%) appliquent un taux comparable dès un
revenu proche du salaire moyen ; l'Allemagne applique 47,5~\% à partir de quatre fois ce salaire
(figure~\ref{fig:seuil}).

\figbloc{r6_taux_superieur_et_seuil}{Taux marginal supérieur de l'impôt sur le revenu et seuil
d'application dans l'OCDE}{fig:seuil}

La trajectoire française du taux supérieur est la plus singulière des grands pays.
Mankiw, Weinzierl et Yagan documentent la baisse générale entre 1983 et 2007 : le taux marginal
supérieur combiné à la TVA est passé dans l'OCDE de près de 80~\% à un peu plus de 60~\%. En
France, le taux statutaire supérieur passe de 58,3~\% en 2000 à 45,8~\% en 2006, puis remonte
jusqu'à 55,4~\% en 2018 et s'y maintient. L'Allemagne baisse de six points et s'arrête, la Suède,
les États-Unis et l'Estonie de trois à quatre, le Royaume-Uni monte de dix points en 2010 puis en
rend cinq (figure~\ref{fig:tauxsuph}).

\figbloc{c21_taux_superieurs}{La France a abandonné douze points de taux supérieur, puis les a
repris}{fig:tauxsuph}

Le prélèvement sur l'euro suivant est le plus lourd au bas de l'échelle. Pour un
célibataire sans enfant, le coin fiscal marginal atteint 64,6~\% à 67~\% du salaire moyen,
redescend à 58,2~\% au salaire moyen, puis remonte à 60,0~\% à 167~\%
(figure~\ref{fig:coinniveau}). Personne n'a voulu ce creux ; il résulte de la superposition des
allègements de cotisations, de la prime d'activité et des aides sous condition de ressources,
dont chacun est défendable isolément.

\figbloc{c4_coin_par_niveau}{Le profil français des prélèvements sur le travail n'est pas
monotone}{fig:coinniveau}

Les effets de ce creux sur la distribution des salaires sont plus difficiles à voir que le creux lui-même. Les plafonds des deux baisses proportionnelles, à 2,5 et 3,5 fois le salaire minimum, étaient
de vrais seuils : un euro de plus faisait perdre six points, ou 1,8 point, de cotisations sur la
totalité du salaire. La Cour des comptes ne trouve pourtant aucune concentration de salariés juste
en dessous de ces plafonds, ni de celui de 1,6 fois le salaire minimum. La part des salariés payés
au salaire minimum, montée à 17,3~\% en 2023 avec les revalorisations, est revenue à 14~\% en 2024.
Le risque d'une «~trappe à bas salaires~» est inscrit dans le barème ; il n'apparaît pas, à ce
jour, dans la distribution observée.

Passer d'un mi-temps à un temps plein coûte d'autant plus cher qu'on est déjà au-dessus
du salaire minimum. La part du salaire supplémentaire absorbée par les cotisations, l'impôt et le
retrait des prestations est de 32~\% au salaire minimum, quel que soit le ménage, mais de 41 à
54~\% à 67~\% du salaire moyen et au salaire moyen. Pour un célibataire à 67~\% du salaire moyen,
elle atteint 51~\%, ce qui place la France au deuxième rang des pays comparés, derrière la
Belgique (figure~\ref{fig:metr}).

\figbloc{m3_taux_marginaux_effectifs}{Ce que coûte une augmentation du temps de travail, par type
de ménage et par pays}{fig:metr}

Le niveau moyen du prélèvement raconte une autre histoire. Le coin fiscal moyen d'un célibataire sans enfant atteint 41,2~\% du coût du
travail à 67~\% du salaire moyen, 47,2~\% au salaire moyen et 54,1~\% à 167~\%
(figure~\ref{fig:cointrois}). L'écart entre les deux extrémités, treize points, est plus faible
qu'en Italie, en Irlande ou aux Pays-Bas. Mais il part d'un niveau que la moyenne de l'OCDE
n'atteint pas même au plus haut des trois points : à 67~\% du salaire moyen, seules la Belgique,
l'Allemagne et l'Autriche prélèvent davantage, et la Hongrie autant. Le problème français n'est
pas la pente, c'est le point de départ --- ce que la figure~\ref{fig:ecarts} montrait déjà, avec
5,7 points de PIB de cotisations sociales au-dessus de la moyenne.

\figbloc{c16_coin_trois_niveaux}{Le coin fiscal français est élevé dès le bas de l'échelle, et le
reste}{fig:cointrois}

\subsection{Taxer sur ce qui ne se triche pas}

Puisque le problème vient de ce que l'État n'observe pas la capacité, tout indice
observable et non manipulable qui lui est corrélé améliore le système. Akerlof (1978) a nommé
cette idée le \emph{tagging}. Un transfert conditionné à une caractéristique que personne ne
change pour des raisons fiscales ne décourage pas le travail, puisqu'il ne diminue pas quand on
gagne davantage. Les prestations liées au handicap, les minima vieillesse et les allocations
familiales reposent sur ce principe : elles atteignent des populations précises sans créer les
taux marginaux élevés qu'engendre toute condition de ressources. La condition de validité est
stricte : il faut que les personnes à forte capacité ne puissent pas se faire passer pour membres
du groupe ciblé, et que le signal soit net.

La théorie va ici beaucoup plus loin que ce qu'aucun pays n'accepte. Elle recommande que l'impôt dépende de toute caractéristique largement exogène, facile à
contrôler et liée à la capacité ; Mankiw, Weinzierl et Yagan en donnent la liste sans euphémisme :
le sexe, la taille, la couleur de peau, l'apparence physique, la santé, l'éducation des parents.
Des travaux sérieux ont chiffré les gains d'un impôt dépendant du sexe ou de la taille, et les ont
trouvés substantiels. Aucun système fiscal moderne ne le fait. Les sociétés acceptent de cibler
sur des caractéristiques qui surviennent au cours d'une vie --- devenir parent, devenir handicapé,
vieillir --- et refusent celles qui sont fixées à la naissance ; la seule exception admise est
l'âge.

\begin{encadre}{Pourquoi cette frontière n'est pas un préjugé}
Un système qui cible une caractéristique pouvant survenir à n'importe qui ressemble à une
assurance : avant de savoir si l'on sera handicapé, chacun a intérêt à ce qu'un dispositif allège
la charge de ceux qui le deviendront. Tout le monde peut y souscrire derrière le voile
d'ignorance.

Un système qui cible une caractéristique connue dès la naissance n'est plus une assurance : c'est
un transfert pur, et ceux qui savent déjà qu'ils sont du «~mauvais~» côté s'y opposent
rationnellement. Deux principes extérieurs au modèle formalisent ce refus : l'équité horizontale
--- deux personnes de capacité semblable doivent payer autant --- et le principe du bénéfice ---
l'impôt devrait être lié à ce que l'on reçoit de la collectivité.
\end{encadre}

\subsection{Ce qu'il faudrait faire}

Lisser le creux au-dessus du salaire minimum avant de toucher au sommet. C'est là que le
prélèvement sur l'euro suivant est le plus lourd et que l'effectif concerné est le plus nombreux.
Le moyen est de reprofiler ensemble la sortie des allègements de cotisations et celle de la prime
d'activité, qui sont aujourd'hui conçues séparément.

La réforme engagée en 2026 va dans ce sens. La loi de financement de la
sécurité sociale pour 2025 fusionne en 2026 les trois dispositifs en une seule réduction dégressive,
plafonnée à trois fois le salaire minimum : les seuils disparaissent, remplacés par une pente plus
longue. Le CAE recommandait en 2019 d'aller plus loin, en concentrant les allègements sous 1,6 fois
le salaire minimum et en abandonnant ceux qui portent au-delà de 2,5, voire de 1,6, si les
évaluations confirmaient leur absence d'effet ; la Cour des comptes propose en outre d'inclure les
compléments de salaire dans la rémunération qui sert à les calculer. La limite est connue : selon une
simulation que cite la Cour, supprimer tous les allègements détruirait de l'ordre d'un million
d'emplois. Il s'agit de reprofiler la pente, pas de la supprimer.

Fusionner les deux impôts sur le revenu. La France prélève 153~Md\euro{} de contribution
sociale généralisée et 96~Md\euro{} d'impôt sur le revenu, sur des assiettes différentes, l'une
sur l'individu, l'autre sur le foyer, l'une proportionnelle, l'autre progressive. Un impôt unique,
de structure lisible, avec un abattement plutôt qu'un empilement de tranches, ferait le même
travail redistributif et rendrait visible ce que chacun paie à la marge.

\subsection{L'objection la plus forte}

On objectera que les cotisations ouvrent des droits et que les fondre dans l'impôt détruirait la logique contributive. C'est le principe du bénéfice, et il est sérieux. Une
cotisation de retraite proportionnelle au salaire, qui ouvre une pension proportionnelle aux
cotisations, est en partie une épargne forcée et non un prélèvement sans contrepartie ; le
salarié la perçoit comme telle, et son effet sur l'offre de travail est plus faible que celui d'un
impôt du même montant.

L'objection ne vaut que pour la part contributive. Les cotisations
d'assurance maladie et de famille n'ouvrent aucun droit proportionnel à ce qui est versé, pas plus
que la contribution sociale généralisée. La comparaison avec le Danemark montre que l'étiquette ne
change rien à ce que le salarié perd sur l'euro suivant. Une réforme peut donc préserver la
logique contributive là où elle existe --- les retraites --- et traiter le reste pour ce qu'il est :
un impôt sur le revenu du travail.

% ===========================================================================
\section{La TVA}

\begin{argument}
La TVA est le meilleur impôt dont dispose un État, à condition d'être prélevée au même taux sur
tout. Chaque taux réduit est une subvention à un type de consommation, et elle donne plus d'euros
aux ménages aisés qu'aux ménages modestes.
\end{argument}

\subsection{L'exemple : un taux réduit sur l'alimentation}

Un taux réduit donne plus d'euros aux riches qu'aux pauvres. Un ménage modeste consacre
une part plus élevée de son revenu à l'alimentation --- disons 18~\% de 1\,500~\euro{} par mois,
soit 270~\euro{} hors taxe. Un ménage aisé y consacre une part plus faible, 9~\% de 5\,000~\euro{},
mais cela fait 450~\euro{}. Quinze points de TVA en moins valent donc 40~\euro{} par mois au
premier et 68 au second (figure~\ref{fig:tauxreduit}). La mesure est progressive en pourcentage du
revenu et régressive en euros. Or ce qui achète du pain, ce sont des euros : verser directement
68~\euro{} au ménage modeste coûterait moins cher que d'en distribuer 40 à l'un et 68 à l'autre.

\figbloc{c3_taux_reduit}{Un taux réduit donne plus d'euros au ménage aisé qu'au ménage
modeste}{fig:tauxreduit}

\subsection{Ce que dit la théorie}

Si l'on dispose d'un barème sur le revenu bien construit, il ne sert à rien d'appliquer
des taux différents aux différents biens. C'est le résultat d'Atkinson et Stiglitz (1976). Deux
personnes qui gagnent le même revenu et paient le même impôt consomment différemment selon leurs
goûts ; taxer plus lourdement le bien qu'aime la première ne redistribue rien, cela pénalise un
goût et non une situation. La formulation exacte contient sa limite : le résultat suppose que la
façon de répartir son budget ne dépend pas du temps passé à travailler, ni de la capacité, de
sorte que le choix de consommation ne révèle rien que le revenu ne révèle déjà.

La théorie admet trois écarts à l'uniformité. Les biens
complémentaires du travail --- garde d'enfants, transport domicile-travail, repas pris dehors
parce qu'on travaille --- peuvent être traités à part. Les biens qui révèlent une capacité que le
revenu déclaré ne montre pas peuvent être taxés davantage. Et si la structure des taxes modifie le
salaire relatif de différents types de travailleurs, l'uniformité peut cesser d'être optimale ---
argument correct en théorie, qu'on ne sait pas mobiliser dans un cas concret. L'architecture
française des taux réduits ne s'explique par aucune de ces trois exceptions, mais par l'histoire
et par les rapports de force sectoriels.

C'est la leçon de la théorie que la pratique a le mieux suivie. Une TVA bien construite ne taxe pas les biens intermédiaires et taxe
uniformément les biens finals : elle applique à la fois Diamond-Mirrlees et Atkinson-Stiglitz.
Vingt-neuf des trente pays de l'OCDE en disposaient en 2009, contre douze en 1976 ; sur ces douze,
onze ont relevé leur taux, de 15,6 à 20,4~\% en moyenne, et la France est le seul à l'avoir
abaissé, de 20 à 19,6~\%. La part des recettes tirées des taxes générales sur la consommation a
doublé entre les années 1960 et le milieu des années 2000, pendant que celle des accises reculait
de 24,1 à 11,5~\%. Sur l'uniformité, en revanche, presque tous les pays ont gardé des taux réduits
pour les «~biens de base~».

\subsection{Ce que montrent les données}

La TVA française a un taux dans la moyenne et une assiette plus étroite que la moyenne.
Son taux normal de 20~\% est proche de la moyenne de l'OCDE (19,4~\%) et inférieur à la médiane
(21~\%). Son ratio de recettes --- ce qu'elle rapporte, rapporté à ce qu'elle rapporterait si toute
la consommation était taxée au taux normal --- est de 51~\%, contre 55~\% en moyenne ; la
Nouvelle-Zélande, qui taxe presque tout à 15~\%, approche 100~\% (figure~\ref{fig:tva}). Ce ratio
mélange les taux réduits, les exonérations de secteurs entiers et la fraude ; on ne peut donc pas
lire l'écart à 100~\% comme un gisement de recettes, puisqu'une partie correspond à des
exonérations que personne ne propose de lever, à commencer par les loyers.

\figbloc{r7_tva_taux_et_assiette}{Taux normal et assiette de la TVA dans l'OCDE}{fig:tva}

Le Conseil des prélèvements obligatoires a chiffré ce que coûtent les exceptions. Son
rapport de février 2023 sur la TVA évalue à au moins 47~Md\euro{} le manque à gagner dû aux taux
réduits et aux autres dérogations au taux normal, 24~\% de ce que l'impôt rapportait en 2021. Une
centaine de catégories de biens et de services en bénéficient, mais les dix premières mesures font
plus de 82~\% du total. Seuls 17~Md\euro{} sont comptés comme dépenses fiscales ; le reste est
considéré comme faisant partie de la norme, et n'apparaît donc pas dans la liste officielle des
niches.

Rapportée au revenu d'une année, la TVA est régressive. Sur données françaises, le CPO la mesure à 12,5~\% du revenu disponible du
dixième le plus modeste et à 4,7~\% de celui du dixième le plus aisé ; l'écart s'explique
entièrement par la part du revenu consommée. C'est l'erreur de dénominateur du chapitre~3 : les
ménages dont le revenu est temporairement bas consomment plus qu'ils ne gagnent, et paraissent donc
lourdement taxés si l'on rapporte leur TVA à ce revenu. Une analyse sur le cycle de vie, qui tient
compte de ce que l'épargne d'aujourd'hui sera consommée plus tard, réduit la régressivité de près de
moitié sans l'annuler. Ce qui reste tient en partie à ce que les ménages aisés ne consomment pas
toutes leurs ressources de leur vivant : ce qu'ils transmettent échappe à la TVA, et le CPO renvoie
ce point à l'impôt sur les successions (chapitre~9). Le constat vaut aussi pour les ménages
durablement modestes, et c'est pourquoi toute réforme doit venir avec sa compensation.

\subsection{Ce qu'il faudrait faire}

Élargir l'assiette plutôt que relever le taux, et compenser les ménages modestes par un
transfert. C'est la recommandation de la \emph{Mirrlees Review}, qui a chiffré un paquet neutre
à la fois sur le rendement et sur la distribution : convergence des taux réduits vers le taux
normal, et retour du produit aux premiers déciles par le barème direct et les prestations. Le
document de l'OCDE arrive à la même conclusion par une autre voie : élargir l'assiette d'une taxe
sur la consommation vaut mieux qu'en relever le taux, parce qu'un taux élevé nourrit l'économie
souterraine. Le taux normal français n'a pas besoin d'être relevé ; ce sont les trois taux réduits
--- 10, 5,5 et 2,1~\% --- qu'il faut rapprocher.

Le CPO a fait le calcul sur les ménages français. On a vu au chapitre~3 que, pour un même coût, un
chèque ciblé relève près de cinq fois plus le niveau de vie des plus modestes qu'une TVA nulle sur
l'alimentation, même en supposant que 80~\% de la baisse de TVA se retrouve dans les prix. Une
structure de taux plus différenciée encore --- 2,1~\% sur toute
l'alimentation, le taux normal sur la restauration --- ne réduirait que d'un demi-point l'écart de
huit points de taux d'effort entre le premier et le dernier dixième. Un taux réduit redistribue mal
parce qu'il ne peut pas distinguer qui achète. Le CPO recommande donc de supprimer les taux réduits
dont l'évaluation confirme l'inefficacité, ou à défaut de les relever d'un cran dans le barème, et
de soutenir le pouvoir d'achat des ménages modestes par des transferts ciblés.

\subsection{L'objection la plus forte}

Les taux réduits soutiennent l'emploi dans des secteurs qui en créent beaucoup, comme la
restauration ou le bâtiment. L'argument est empirique, et il a été testé. La baisse de la TVA
sur la restauration, passée de 19,6 à 5,5~\% en 2009, a été évaluée par Benzarti et Carloni
(\emph{American Economic Journal: Economic Policy}, 2019), qui comparent les restaurants à d'autres
services marchands sur données d'entreprises. Les propriétaires ont capté 41~\% de l'avantage sous
forme de profits, les salariés 25~\%, les fournisseurs 16~\%, et les consommateurs 18~\%
seulement, parce que les prix ont à peine baissé. Les relèvements ultérieurs, à 7~\% en 2012 puis
10~\% en 2014, sont davantage passés dans les prix que la baisse ne l'avait fait. Le test du statu
quo du chapitre~3 dit le reste : un taux réduit est une subvention à un secteur, financée par un
taux plus élevé sur tous les autres, et il doit se justifier comme telle.

% ===========================================================================
\section{L'épargne et le capital des ménages}

\begin{argument}
Taxer le rendement de l'épargne revient à taxer davantage celui qui consomme plus tard que celui
qui consomme tout de suite. La théorie ne recommande pas de ne pas taxer le capital ; elle
recommande de taxer la rente plutôt que la patience.
\end{argument}

\subsection{L'exemple : deux salariés identiques}

Deux personnes gagnent le même salaire et paient le même impôt dessus. Si le rendement de l'épargne est taxé, celle qui attend pour consommer paie un second impôt que l'autre ne paie pas. À
3~\% de rendement réel et 30~\% de prélèvement, celle qui épargne trente ans avant de consommer
subit l'équivalent d'un impôt de 23~\% sur sa consommation, contre zéro pour l'autre ; à quarante
ans, 30~\% (figure~\ref{fig:attente}). La seule différence entre elles est le moment où elles
consomment.

\figbloc{p4_impot_sur_l_attente}{L'impôt implicite sur la consommation différée croît avec la
durée}{fig:attente}

C'est la violation la plus extrême du principe d'uniformité du chapitre~7. Consommer en
2050 et consommer en 2026 sont deux biens différents, et un impôt sur le rendement les taxe à des
taux différents, d'autant plus que l'échéance est lointaine. Si l'on juge injustifié d'appliquer
5,5~\% aux pâtes et 20~\% aux biscuits, il faut expliquer pourquoi on applique 0~\% à la
consommation de cette année et l'équivalent de 30~\% à celle qu'on diffère de quarante ans.

\subsection{Le mécanisme : rendement normal et rente}

Le rendement normal rémunère le seul fait d'attendre ; la rente est tout ce qui dépasse.
Taxer le premier taxe la patience, et rien d'autre. La rente --- une position de marché, une
information que les autres n'avaient pas, un terrain bien placé --- serait apparue de toute façon :
la taxer ne décourage rien. C'est la distinction la plus utile de toute la fiscalité du capital,
et le système français ne la fait nulle part.

Le mécanisme qui la met en \oe uvre est plus simple que son nom. On laisse le
contribuable déduire chaque année, du revenu de son placement, un montant égal au capital investi
multiplié par un taux sans risque --- celui de la dette publique, par exemple. Ce qui dépasse est
imposé au barème ordinaire ; ce qui reste en deçà ne l'est pas, et le solde non utilisé est
reporté avec intérêts. Le placement qui ne rapporte que le rendement normal n'est jamais imposé ;
la rente l'est intégralement au barème, donc souvent plus lourdement qu'aujourd'hui. La Norvège
applique ce mécanisme aux dividendes et plus-values des actionnaires depuis 2006 ; la
\emph{Mirrlees Review} le recommande pour toute l'épargne, avec un avantage supplémentaire pour
l'épargne retraite.

Le résultat classique selon lequel l'impôt optimal sur le capital serait nul à long terme
est beaucoup moins solide qu'on ne le dit. Les modèles de Chamley et Judd concluaient en ce sens
parce que, si les ménages ont un horizon infini, leur épargne devient à long terme parfaitement
sensible au rendement net : un impôt sur le capital ne réduit pas ce rendement net, il élève le
rendement brut exigé, donc contracte le stock de capital et la production. Straub et Werning
(2020) ont montré que la conclusion dépend de paramètres pour lesquels elle s'inverse. Si les
horizons sont courts, ou si les ménages épargnent par précaution faute de pouvoir s'assurer, taxer
le capital peut redevenir souhaitable. Aucune de ces objections ne justifie de taxer le rendement
normal plus lourdement à mesure qu'on attend : elles portent sur le niveau, pas sur l'uniformité.

Dans les modèles dynamiques, l'impôt sur le capital sert d'abord d'instrument d'information, un résultat que rappellent Mankiw, Weinzierl et Yagan. Le problème de
la redistribution est qu'elle incite à travailler moins pour être mieux traité ; dans une économie
dynamique, on peut en outre épargner pour maintenir sa consommation les années où l'on gagne peu.
Un impôt sur le capital plus lourd quand le revenu du travail chute rend cette stratégie coûteuse.
Kocherlakota (2005) montre que l'espérance de cet impôt, calculée avant que chacun connaisse sa
capacité, est nulle : il ne lève, en moyenne, aucune recette. Aucun conseiller ne proposerait un
impôt de ce genre à un ministre, et les auteurs le disent. Le résultat montre seulement que la
question utile, pour l'impôt sur le capital, est «~pour quoi faire~».

\subsection{Ce que montrent les données}

Le système taxe le rendement nominal, donc taxe l'inflation sans que personne l'ait
voté. À 2~\% d'inflation et 3~\% de rendement réel, un prélèvement de 30~\% porte sur 5~\%, soit
la moitié du rendement réel : le taux affiché est 30~\%, le taux effectif 50~\%. L'impôt absorbe la
totalité du gain réel à 7~\% d'inflation pour un rendement réel de 3~\%, et dès 4,7~\% pour un
rendement réel de 2~\% (figure~\ref{fig:inflation}). Le défaut est connu et reste non corrigé parce
qu'il rapporte, et que ce qu'il rapporte n'apparaît dans aucune ligne de la loi de finances.

\figbloc{c5_inflation}{Un impôt assis sur le rendement nominal taxe l'inflation}{fig:inflation}

Le même rendement réel supporte en France de zéro à 97~\% d'impôt, selon la seule
enveloppe qui le porte. Un livret réglementé est exonéré ; un plan d'épargne en actions après cinq
ans supporte 17,2~\% affichés, soit 28,7~\% sur le rendement réel ; le prélèvement forfaitaire
unique de 2025, 30~\%, en prend 50~\%, et 52~\% à son taux de 2026, 31,4~\% ; un revenu foncier imposé au taux marginal de 41~\% en prend 97~\%
(figure~\ref{fig:supports}). Aucun de ces écarts ne correspond à une différence économique entre
les placements. L'épargnant place donc là où le rendement après impôt est le plus élevé, pas là où
son argent serait le plus utile, et le coût consiste en projets non financés dont on ne connaît ni
le nom ni le montant.

\figbloc{c18_supports_epargne}{Le même rendement réel de 3~\% supporte de 0 à 97~\% d'impôt selon
le support}{fig:supports}

Au sommet, le revenu de l'épargne en actions est taxé en France nettement au-dessus de la
moyenne. Le taux personnel supérieur est de 34~\% sur les dividendes comme sur les plus-values,
contre 25 et 20~\% en moyenne dans l'OCDE, ce qui place la France au huitième et au sixième rang
(figure~\ref{fig:dividendes}). Pour les dividendes, l'impôt sur les sociétés s'y ajoute : c'est la
cascade du chapitre~5.

\figbloc{r8_dividendes_plus_values}{Taux personnels supérieurs sur les dividendes et les
plus-values dans l'OCDE}{fig:dividendes}

\subsection{Taxer le stock plutôt que le flux}

Un impôt annuel sur un patrimoine ne se compare à un impôt sur le revenu qu'une fois converti en taux sur le rendement. Un prélèvement de 1~\% par an sur un capital
qui rapporte 2~\% réels absorbe la moitié du rendement ; le même prélèvement sur un capital qui
rapporte 5~\% en absorbe un cinquième (figure~\ref{fig:stockflux}). Le taux affiché est identique,
le taux réel ne l'est pas, et un impôt sur la fortune frappe donc le plus lourdement les actifs qui
rapportent le moins.

\figbloc{p3_stock_contre_flux}{Part du rendement réel absorbée par un impôt annuel sur le
patrimoine}{fig:stockflux}

La proposition d'impôt minimum sur les très grands patrimoines se juge avec cette arithmétique. Elle consiste à prélever chaque année 2~\% des
patrimoines supérieurs à 100~millions d'euros, l'impôt sur le revenu déjà acquitté venant en
déduction. Son argument le plus fort est réel : une large part du revenu économique des très
grandes fortunes reste logée dans des sociétés holding, où les dividendes reçus des filiales sont
exonérés à 95~\%, de sorte que le revenu imposable du détenteur peut être très inférieur à ce que
son patrimoine lui rapporte. L'impôt minimum ne mord que dans ce cas, et c'est sa justification.

Son défaut apparaît dès la conversion : 2~\% du patrimoine, c'est 67~\% d'un rendement réel de 3~\%, et 100~\% d'un rendement de 2~\%. Pour le fondateur d'une entreprise non cotée valorisée
un milliard, cela représente 20~millions d'euros à trouver chaque année, qu'elle distribue ou non :
il faut alors céder des titres ou emprunter, et la valeur retenue pour une société non cotée est
contestable par construction. L'impôt frappe en outre le rendement normal autant que la rente, ce
qui est exactement ce que la théorie déconseille. L'argument qui le fonde appelle donc un
instrument plus précis : imposer directement le revenu accumulé dans les holdings patrimoniales au
taux de leur détenteur, ce qui atteint les mêmes personnes sans taxer l'attente.

\subsection{Ce qu'il faudrait faire}

Un seul traitement pour toute l'épargne : abattement pour rendement normal, imposition de
ce qui dépasse au barème, assiette indexée sur l'inflation. Les enveloppes perdraient leur raison
d'être, puisque le rendement normal ne serait plus taxé nulle part ; la rente le serait partout, au
même taux. L'épargne retraite peut garder un avantage, que la \emph{Mirrlees Review} justifie non par la
neutralité, mais par le manque de prévoyance de certains épargnants et par ce que coûtent à la
collectivité ceux qui arrivent à la retraite sans épargne.

\subsection{L'objection la plus forte}

Le patrimoine est très concentré, et exonérer le rendement normal profiterait d'abord aux
plus riches. L'objection est sérieuse, puisque la part du rendement normal dans le revenu des
patrimoines importants est élevée en valeur absolue.

Elle confond pourtant deux questions. Le rendement normal d'un grand patrimoine reste
taxé au moment de sa transmission, dont le chapitre~9 montre qu'elle rapporte déjà davantage en
France que partout ailleurs dans l'OCDE. La rente, elle, serait imposée au barème, donc, pour les
hauts revenus, à un taux supérieur au prélèvement forfaitaire actuel. Et la concentration du patrimoine appelle un impôt sur les
transmissions bien construit, pas un impôt sur l'attente qui pèse sur tous les épargnants, y
compris les plus modestes dont l'épargne dort sur des supports peu rémunérés.

% ===========================================================================
\section{L'immobilier et les transmissions}

\begin{argument}
La valeur du sol est la rente la plus pure qui existe : la taxer ne décourage rien, puisque
personne ne produit de terrain. Taxer la transaction est à l'inverse la pire construction
possible, puisqu'elle frappe celui qui bouge et laisse intact celui qui ne bouge pas.
\end{argument}

\subsection{L'exemple : un impôt annoncé est payé d'un coup}

Un impôt sur la valeur du sol ne peut pas en réduire la quantité ; c'est son prix qui
baisse. La valeur d'un terrain est le loyer qu'il rapportera, actualisé. Un impôt annuel de 1~\%
sur un terrain actualisé à 4~\% fait tomber sa valeur de 20~\% le jour de l'annonce ; un impôt de
2~\%, de 33~\% (figure~\ref{fig:capitalisation}). Tous les acquéreurs suivants achètent moins cher
et ne supportent rien. C'est ce qui rend l'impôt efficient --- il ne pèse sur aucune décision future
--- et c'est ce qui le rend politiquement brutal : une seule génération de propriétaires acquitte
la totalité.

\figbloc{c19_capitalisation}{Un impôt foncier annoncé est payé d'un coup par le propriétaire du
jour}{fig:capitalisation}

Un droit de mutation, lui, ne frappe que ceux qui déménagent, et il décourage le déménagement.
Au Royaume-Uni, le passage du droit de 1 à 3~\% au-delà de 250\,000~£ réduit d'environ 37~\% la
mobilité des ménages concernés (Hilber et Lyytikäinen, 2017), et la suppression temporaire d'un
droit de 1~\% en 2008-2009 a augmenté les transactions de 20~\% (Best et Kleven, 2018). L'effet
porte surtout sur les déménagements courts, liés au logement ; Hilber et Lyytikäinen ne trouvent
pas d'effet sur les mobilités liées à l'emploi, si bien que le coût pour l'appariement entre
travailleurs et emplois, souvent invoqué, est moins établi. En France, la hausse de 0,7 point des
droits en 2014 n'a réduit les transactions qu'à court terme (Bérard et Trannoy, 2018). Le droit
laisse intacte la situation de celui qui ne bouge pas, quelle que soit la valeur de son bien, et il
est payé par le vendeur plus que par l'acheteur qui le verse (chapitre~2).

\subsection{Ce que disent les trois textes}

La \emph{Mirrlees Review} recommande d'abandonner les taxes sur les transactions au profit d'un impôt sur la valeur du sol. Elle va plus loin pour le logement occupé par son propriétaire. Deux
ménages de même revenu, dont l'un possède son logement et l'autre loue le sien, ne sont pas traités
de la même façon : le propriétaire consomme un service de logement que rien ne taxe, tandis que le
locataire paie un loyer sur un revenu déjà imposé. C'est le raisonnement du chapitre~7 --- deux
formes de la même consommation taxées différemment --- et sa conclusion, l'imposition d'un loyer
implicite, est électoralement intenable dans la plupart des pays où elle a été énoncée. Les
Pays-Bas, qui imposent un loyer implicite mais laissent déduire les intérêts d'emprunt, ont
recommandé en 2024 d'éteindre l'avantage net qui en résulte.

Le document de l'OCDE place les impôts récurrents sur l'immobilier au premier rang des
impôts les moins défavorables à la croissance. C'est le résultat le plus repris de ce travail, et
c'est aussi celui dont la robustesse a été la plus contestée (chapitre~11). Il coïncide en tout cas
avec ce que la théorie prédit pour des raisons indépendantes : une assiette immobile se taxe sans
dommage.

\subsection{Ce que montrent les données}

La France taxe deux fois plus que la moyenne de l'OCDE à la fois la détention et la
transaction. Les impôts récurrents sur la propriété immobilière atteignent 1,9~\% du PIB, contre
0,95~\% en moyenne ; les droits sur les transactions, 0,7~\% contre 0,37~\% (figure~\ref{fig:immo}).
Le partage entre les deux est donc celui de la moyenne. Ce qui distingue la France est le niveau
des droits sur les transactions et l'assiette de l'impôt récurrent.

\figbloc{r9_detention_et_transaction}{Impôts récurrents sur l'immobilier et droits sur les
transactions dans l'OCDE}{fig:immo}

Les droits de mutation suivent le marché, ce qui en fait aussi une recette instable. Ils
sont tombés de 0,74 à 0,50~\% du PIB entre 2022 et 2024, alors que la taxe foncière n'a pas bougé
(figure~\ref{fig:detention}). Pour les départements, qui en tirent l'essentiel, c'est une recette
qui baisse exactement quand l'activité ralentit. En 2024, les droits d'enregistrement sur les
mutations à titre onéreux ont rapporté 14,7~Md\euro{} ; les taxes foncières et l'impôt sur la
fortune immobilière, 45,6~Md\euro{}.

\figbloc{g10_detention_transaction}{Taxer la détention, la transmission ou la transaction, en \%
du PIB}{fig:detention}

L'assiette de l'impôt récurrent est fausse. La
taxe foncière des logements est assise sur des valeurs locatives établies en 1970 et revalorisées
depuis par coefficients forfaitaires. Deux logements de même valeur de marché peuvent donc supporter des
impôts très différents selon que leur quartier s'est valorisé ou déprécié depuis cinquante-cinq
ans. Tant que l'assiette reste fausse, toute hausse de l'impôt foncier amplifierait des écarts
injustifiés.

Cette assiette rend aussi l'impôt régressif. Une étude de l'Insee pour le Conseil des prélèvements obligatoires, publiée avec son
rapport de décembre 2023 sur le logement, montre que, parmi les propriétaires, la taxe foncière
pèse relativement plus sur les moins aisés. La raison principale est l'assiette : les valeurs de
1970 sous-évaluent davantage les logements des communes les plus aisées. S'y ajoute que les
patrimoines immobiliers les plus importants se trouvent dans des communes aux taux plus bas. Le
même rapport chiffre la fiscalité du logement à 92~Md\euro{} en 2022, près de
8~\% des prélèvements obligatoires, dont 15~Md\euro{} de dépenses fiscales réparties entre
soixante-dix dispositifs, pour la plupart ni bornés dans le temps ni évalués.

\subsection{Les transmissions}

Aucun pays de l'OCDE ne prélève autant que la France sur les successions et les
donations. Ces impôts y rapportent 0,76~\% du PIB, cinq fois la moyenne de l'OCDE (0,15~\%) ; la
Corée, la Belgique et le Japon suivent autour de 0,6~\% (figure~\ref{fig:successions}).

\figbloc{r10_successions}{Impôts sur les successions et donations dans l'OCDE}{fig:successions}

Ce rendement élevé est obtenu avec un barème élevé sur une assiette percée. Le barème en
ligne directe monte jusqu'à 45~\%, mais l'abattement de 100\,000~\euro{} par parent et par enfant
se reconstitue tous les quinze ans, l'assurance vie suit un régime propre, et les transmissions d'entreprises
bénéficient d'une exonération de 75~\% de leur valeur sous condition d'engagement de conservation.
Le résultat est l'inverse de la règle du chapitre~1 : des taux élevés sur une assiette étroite,
donc une forte incitation à organiser sa transmission, et une charge qui dépend moins du montant
transmis que de la qualité du conseil.

Le Conseil d'analyse économique a mesuré ce que l'assiette percée fait à la
progressivité. Sa note de décembre 2021, signée par Clément Dherbécourt, Gabrielle Fack, Camille
Landais et Stefanie Stantcheva, s'appuie sur les données de transmission de la DGFiP, complétées
par les comptes de patrimoine. Le millième le plus aisé de chaque génération, qui reçoit au cours
de sa vie de l'ordre de 13~millions d'euros, acquitte à peine 10~\% de droits sur l'ensemble, loin
du taux marginal de 45~\% applicable au-delà de 1,8~million en ligne directe. Entre le
quatre-vingt-dixième centile et ce millième, les sommes héritées vont de 1 à 24, les taux effectifs
de 1 à 2 : la charge est relativement lourde pour le bas du dernier décile des héritages, ce qui
nourrit l'idée que l'impôt frappe l'épargne d'une vie de travail. Les trois grandes exceptions ---
assurance vie, démembrement de propriété, transmission d'entreprises --- coûtaient selon la note
de 2 à 5~milliards d'euros par an chacune. L'enjeu grandit : l'héritage représente 60~\% du
patrimoine total, contre 35~\% au début des années 1970.

La Cour des comptes a depuis évalué l'une de ces exceptions, le pacte Dutreil. Celui-ci
exonère 75~\% de la valeur d'une entreprise transmise si la famille s'engage à en conserver
le contrôle quelques années ; ajouté à la réduction de moitié accordée aux donations avant 70 ans,
il ramène le taux maximal de 45 à 5,6~\%. Dans son rapport de novembre 2025, la Cour estime qu'il a
coûté 5,5~Md\euro{} en 2024, contre 1,2~Md\euro{} en 2020 --- en partie à cause d'une très grosse
donation en 2023 et en 2024 --- alors que les documents budgétaires l'ont inscrit pour un
demi-milliard par an de 2011 à 2024. Près des deux tiers de l'avantage, 65~\%, vont à 1~\% des
bénéficiaires, 110~personnes en 2024, qui en retirent en moyenne 30~millions d'euros chacune ; l'industrie,
souvent invoquée pour le justifier, en reçoit 17~\%.

L'évaluation mesure aussi ce que le dispositif produit. Avec
l'Institut des politiques publiques, la Cour compare des entreprises transmises entre 2010 et 2018
sous pacte et hors pacte : les premières sont plus grandes en moyenne, mais leur rentabilité et
leur investissement étaient proches avant la transmission. Elles changent moins souvent de
contrôle dans les cinq ans --- un peu moins
de 30~\% contre un peu plus de 40~\% --- et disparaissent un peu moins souvent, 6~\% contre 10~\% au
bout de neuf ans. Elles n'investissent pas davantage et n'emploient pas davantage. Les familles qui
signent un pacte ne sont pas tirées au sort, et la comparaison ne vaut pas expérience ; mais le
dispositif obtient ce qu'il promet sur le contrôle familial, pas ce qu'on en attend sur l'activité.
La Cour admet sa prémisse --- un taux marginal de 45~\% est élevé pour des actifs professionnels
en comparaison européenne --- mais propose de réduire l'abattement, d'en exclure les actifs non
professionnels et la trésorerie excédentaire qui en profitent aujourd'hui, de le rendre dégressif
au-delà d'un certain montant transmis, et de faciliter en contrepartie le paiement différé et
fractionné des droits. La loi de finances pour 2026 n'a retenu qu'une partie de ces pistes : pour
les transmissions réalisées depuis le 21~février 2026, l'engagement individuel de conservation
passe de quatre à six ans, et certains biens --- logements, yachts, aéronefs, objets d'art, vins,
chevaux de course --- sortent de l'exonération lorsque la société ne les affecte pas exclusivement
à son activité depuis au moins trois ans. Le taux de 75~\% n'a pas changé.

La \emph{Mirrlees Review} propose un impôt sur l'ensemble des sommes reçues au cours de
la vie. Toutes les transmissions, entre vifs ou au décès, quelle que soit la forme de l'actif,
seraient additionnées pour chaque bénéficiaire et taxées selon un barème unique. L'argument est
d'équité --- réduire l'inégalité des chances qui naît du hasard de la naissance --- pour un coût
économique faible, puisque l'impôt porte sur ce que le bénéficiaire n'a pas produit. C'est aussi
la proposition du CAE, dont les simulations montrent qu'une assiette élargie et des taux plus bas
pourraient réduire les droits de 99~\% de la population tout en dégageant un surplus de recettes.

\subsection{Ce qu'il faudrait faire}

Réviser les valeurs locatives d'abord, puis déplacer progressivement la charge des droits
de mutation vers l'impôt récurrent, en l'asseyant autant que possible sur le sol. L'ordre importe :
alourdir un impôt assis sur des valeurs fausses serait injuste, et l'asseoir sur le bâti plutôt que
sur le sol revient à taxer le capital investi dans la construction. C'est l'ordre que recommande
le Conseil des prélèvements obligatoires en 2023 : relier d'abord l'assiette de la taxe foncière aux
loyers ou aux prix de marché, par une méthode qu'on puisse actualiser à faible coût, puis réduire
les droits de mutation en relevant l'impôt sur la détention, sans perte pour les collectivités. Il
propose aussi de remplacer les abattements pour durée de détention sur les plus-values
immobilières, qui récompensent l'attente et incitent à garder des terrains inutilisés, par la
seule prise en compte de l'érosion monétaire et des travaux --- la logique d'indexation du
chapitre~8 --- et d'unifier à terme les régimes de la location meublée et de la location nue.

Pour les transmissions : élargir l'assiette, en plafonnant les régimes dérogatoires, et
baisser les taux, à rendement constant. Le CAE en 2021 et la Cour des comptes en 2025 proposent
tous deux de réduire fortement l'abattement Dutreil en facilitant l'étalement du paiement des
droits ; le CAE jugeait même préférable de le supprimer, et chiffrait à au moins 1,5~Md\euro{} le
seul passage de 75 à 50~\% assorti de la fin de la réduction pour donation avant 70 ans. Il y
ajoutait l'intégration de l'assurance vie au barème commun et la taxation de l'usufruit lorsqu'il
rejoint la nue-propriété au décès ; avec l'ensemble de ces mesures et le barème actuel, le
supplément de recettes atteignait 9~Md\euro{}, de quoi baisser nettement les taux à rendement
constant.

\subsection{L'objection la plus forte}

«~C'est un impôt sans distorsion~» et «~c'est une spoliation d'une génération de propriétaires~» sont deux descriptions exactes de la même propriété. La
capitalisation fait payer en une fois les propriétaires du jour de l'annonce, y compris ceux qui
ont un patrimoine immobilier important et des revenus faibles --- souvent des retraités.

Un projet sérieux doit donc traiter cette génération. Une montée en
charge assez lente pour que la capitalisation se diffuse, et la possibilité de reporter le paiement
jusqu'à la vente ou à la succession pour les propriétaires à faible revenu, avec un intérêt modéré.
Sans ces dispositions, la proposition est économiquement juste et politiquement morte, ce qui
revient à ne rien proposer.

% ===========================================================================
\section{Le carbone}

\begin{argument}
Une tonne de CO$_2$ cause le même dommage d'où qu'elle vienne. La taxer à des prix différents
selon le secteur revient à payer plus cher une réduction donnée des émissions, quel que soit le
niveau d'ambition retenu.
\end{argument}

\subsection{L'exemple : deux entreprises, deux tonnes}

Supposons qu'on veuille réduire les émissions de deux tonnes, et que deux entreprises
puissent le faire. La première peut supprimer chaque tonne pour 30~\euro{}, la seconde pour
90~\euro{}. Avec un prix unique de 50~\euro{} la tonne, la première réduit ses deux tonnes, puisque
cela lui coûte moins que de payer, et la seconde ne fait rien : l'objectif est atteint pour
60~\euro{}. Si l'on taxe la seconde à 100~\euro{} et pas la première, c'est la seconde qui réduit
ses deux tonnes : même résultat pour le climat, 180~\euro{} de coût. Le prix unique garantit que les
réductions les moins chères sont faites en premier, où qu'elles se trouvent.

C'est un gain qui ne demande aucun accord préalable sur l'ambition climatique. Même celui
qui juge l'objectif trop élevé a intérêt à ce qu'il soit atteint au moindre coût, puisque ce sont
autant de ressources qui restent disponibles pour autre chose.

\subsection{Ce que montrent les données}

En France, la tonne de CO$_2$ émise sur la route paie plus de 120~\euro{} ; celle émise
dans l'industrie ou les bâtiments, presque toujours moins de 60~\euro{}, quand elle paie quelque
chose. En 2021, 94~\% des émissions du transport routier étaient tarifées au-dessus de
120~\euro{} la tonne ; 41~\% de celles de l'industrie, 36~\% de celles des bâtiments et la moitié
de celles des transports hors route ne supportaient aucun prix (figure~\ref{fig:carbone}). Au
total, un tiers des
émissions de CO$_2$ issues de l'énergie sont tarifées au-dessus de 60~\euro{} la tonne, ce qui
place la France au milieu des pays comparés (figure~\ref{fig:carbone_pays}).

\figbloc{m1_prix_carbone_secteurs}{La dispersion du prix implicite du carbone entre secteurs
français, 2021}{fig:carbone}

\figbloc{m2_carbone_comparaison}{Part des émissions tarifées au-dessus de 60~\euro{} la tonne,
2021}{fig:carbone_pays}

\subsection{Ce qu'il faudrait faire}

Étendre le prix, pas le relever. Le transport routier est le seul secteur où la France
applique déjà un prix élevé et cohérent avec ses objectifs ; la réforme consiste à aligner le reste
sur lui. Les quotas échangeables produisent le même signal qu'une taxe, à une condition que la
\emph{Mirrlees Review} rappelle : les vendre aux enchères plutôt que les distribuer. Un quota
gratuit ne rapporte rien à l'État, donc ne permet d'alléger aucun autre prélèvement, et revient à
offrir une rente aux émetteurs installés.

Les exceptions ont un coût, que la note du CAE de mars 2019 a chiffré. Signée par Dominique
Bureau, Fanny Henriet et Katheline Schubert, elle relevait que 15~\% des
émissions liées à l'énergie échappaient alors à tout prix, ni composante carbone ni quotas
européens. Les principales exonérations étaient le carburant de l'aviation commerciale
(3,6~Md\euro{} en 2019), le remboursement partiel aux transporteurs routiers (1,5~Md\euro{}) et le
taux réduit des exploitants agricoles (0,9~Md\euro{}). Une partie des écarts de taxe sur les
carburants se défend : cette taxe tient aussi lieu de redevance pour l'usage des routes, que les
engins agricoles et de chantier n'empruntent pas --- ce qui n'est pas le cas des camions. La note en
tire une construction simple : faire payer la composante carbone sur toute tonne émise, quel qu'en
soit l'usage, et ne moduler que le reste.

\subsection{L'objection la plus forte}

Un prix du carbone élevé et uniforme pèse lourdement sur les ménages qui dépendent de la
voiture et sur les industries exposées à la concurrence internationale. La France l'a éprouvé :
la trajectoire de la taxe carbone a été gelée fin 2018 après le mouvement des gilets jaunes. Et une
industrie taxée seule peut perdre des parts de marché au profit de concurrents qui émettent davantage
ailleurs.

La note du CAE a chiffré la première moitié de l'objection. La hausse prévue avant le
gel --- de 44,6 à 86,2~\euro{} la tonne, avec le rattrapage du gazole sur l'essence --- aurait
coûté près de 1~\% de leur revenu disponible aux 10~\% de ménages les plus modestes, contre 0,3~\%
aux plus aisés. Reversée en totalité sous la forme d'un même montant par unité de consommation, la
recette aurait fait gagner en moyenne 60~\euro{} par an au premier décile et perdre 80~\euro{} au
dernier, puisque les ménages modestes émettent moins. Les écarts à l'intérieur d'un même décile
restent pourtant considérables : même après ce reversement, un ménage modeste sur dix perd plus de
80~\euro{} par unité de consommation. Ces écarts tiennent moins au lieu de résidence qu'à
l'équipement : avant reversement et à revenu égal, un ménage rural perd 130~\euro{} de plus par
unité de consommation qu'un ménage de l'agglomération parisienne, mais 20~\euro{} seulement à
chauffage et véhicule identiques.

Les deux problèmes ont des réponses qui préservent le prix. Le premier se traite par un
reversement du produit, qui laisse intact le signal sur chaque litre consommé tout en compensant
les ménages modestes --- c'est la règle du chapitre~3 : juger la réforme avec sa compensation. Un
transfert décroissant avec le revenu et modulé selon le territoire permettait, selon le CAE, de
rendre gagnants plus de 90~\% des ménages des cinq ou six premiers déciles ; ce qui reste relève
d'aides temporaires au remplacement de la chaudière ou du véhicule, pas d'une exonération. Le
second se traite par un ajustement aux frontières, que l'Union européenne applique à plusieurs
industries depuis 2026. Exonérer un secteur, en revanche, revient à payer plus cher chaque tonne
évitée ailleurs.

\partie{Troisième partie}{Le système d'ensemble}

% ===========================================================================
\section{Le niveau des prélèvements et le classement de l'OCDE}

Le problème français tient d'abord à la dispersion des prélèvements. À
43,5~\% du PIB en 2024, cotisations imputées exclues, la France arrive au deuxième rang européen,
à égalité avec l'Autriche et derrière le Danemark, qui est à 45,4~\% sans que personne y voie une
pathologie (figure~\ref{fig:niveau}). Ce qui distingue la France est que, pour un même euro de
richesse, le taux appliqué varie énormément selon la forme que prend cette richesse, le secteur qui
la produit et le support qui la porte. C'est exactement ce que la règle du carré du taux interdit
(chapitre~1). La mesure de l'OCDE, utilisée en tête de la deuxième partie, donne un niveau voisin
--- 43,9~\% en 2023 --- avec des conventions un peu différentes. La figure~\ref{fig:carte} donne la
taille de chaque prélèvement français : cinq d'entre eux --- cotisations employeurs et salariés,
TVA, contribution sociale généralisée, impôt sur le revenu --- font plus des deux tiers du total.

\figbloc{s4_taux_global}{Impôts et cotisations sociales en \% du PIB, 1995-2024}{fig:niveau}

\figbloc{s1_carte_des_prelevements}{Les prélèvements obligatoires français de 2024, un rectangle
par prélèvement}{fig:carte}

\subsection{Le classement de l'OCDE appliqué à la France}

Le résultat qui a fait la fortune du document de l'OCDE tient en une phrase : les impôts
récurrents sur l'immobilier seraient les moins défavorables à la croissance, devant les taxes sur
la consommation et les autres impôts sur le patrimoine, puis les impôts sur le revenu des
personnes, les impôts sur les sociétés étant les plus nocifs. La recommandation qui en découle est
attendue : à rendement constant, déplacer une partie de la charge des impôts sur le revenu vers les
prélèvements moins distorsifs.

Appliqué à la France, le classement dessine une structure presque à l'envers de ce qu'il
recommande, avec une exception. La France prélève 25,9 points de PIB sur le revenu des personnes,
les cotisations et la masse salariale, contre 17,5 en moyenne : c'est là que se concentre l'écart, et
c'est l'avant-dernière catégorie du classement. Elle prélève aussi davantage que la moyenne dans les
catégories les moins nocives --- immobilier, consommation, autres impôts sur le patrimoine. L'exception
est la catégorie jugée la plus nocive, l'impôt sur les sociétés, où elle prélève moins que la moyenne
(figure~\ref{fig:ocde}).

\figbloc{r12_classement_ocde}{La structure française rangée selon le classement de
l'OCDE}{fig:ocde}

\subsection{Ce que vaut ce classement}

La méthode détermine ce que le résultat peut signifier. Les auteurs régressent le
logarithme du PIB par habitant de vingt-et-un pays, entre 1970 et 2005, sur les déterminants
habituels --- capital physique, capital humain, croissance démographique --- en ajoutant le niveau
global des prélèvements, puis la part de chaque catégorie d'impôt dans les recettes. Le niveau étant
contrôlé, les coefficients sur les parts s'interprètent comme l'effet d'un déplacement à rendement
constant. Deux traits de cette construction limitent ce qu'on peut en tirer : le modèle suppose que
les parts et le niveau sont indépendants dans les données, ce que les auteurs eux-mêmes jugent
improbable ; et la comparaison entre pays sur longue période mesure une association entre structures
et niveaux de revenu durablement différents, non l'effet d'une réforme.

Les citations omettent en général l'ampleur des effets estimés, que les auteurs eux-mêmes désavouent. L'estimation centrale est qu'un déplacement de 1~\% des recettes des
impôts sur le revenu vers la consommation et le patrimoine augmenterait le PIB par habitant d'un quart
de point à un point à long terme. Les auteurs écrivent aussitôt : «~L'ordre de grandeur de l'effet
estimé est plus élevé que ce à quoi on pourrait raisonnablement s'attendre~», puis : «~il est clair
que l'ampleur des effets ne peut pas être mesurée avec précision dans un cadre de comparaison entre
pays~».

\begin{encadre}{Pourquoi un quart de point est déjà beaucoup trop}
Un déplacement de 1~\% des recettes représente, en France, environ 13~Md\euro{}. L'estimation basse
du document attribue à ce déplacement un gain permanent d'un quart de point de PIB, soit de l'ordre
de 7~Md\euro{} par an.

Déplacer treize milliards d'une catégorie d'impôt à une autre produirait donc chaque année,
indéfiniment, un gain égal à la moitié du montant déplacé. Aucun mécanisme connu ne produit un
rendement pareil. On peut retenir le signe et l'ordre du classement ; on ne peut pas retenir les
nombres, et chiffrer une réforme française à partir de ces coefficients produirait un résultat
fantaisiste.
\end{encadre}

La robustesse de l'ordre lui-même a été contestée. Xing (2012) montre que le classement ne
résiste pas à tous les changements de spécification, et que l'avantage attribué aux impôts récurrents
sur l'immobilier --- le résultat le plus repris --- disparaît en particulier. Ce qui subsiste n'est
pas rien : l'ordre obtenu coïncide avec ce que la théorie prédit pour des raisons entièrement
indépendantes, puisqu'une assiette immobile se taxe sans dommage et qu'un prélèvement sur un facteur
de production fait le plus de dégâts. Le classement est donc une confirmation plausible d'un
mécanisme établi ailleurs, non une preuve autonome, et c'est sous ce statut qu'il est cité dans ce
cours.

Les recommandations les plus solides du document ne sont d'ailleurs pas dans le classement.
L'essentiel des gains, écrivent les auteurs, vient de la réforme de chaque impôt, pas du déplacement
entre impôts : élargir les assiettes et baisser les taux. Ces recommandations ont été reprises dans
les chapitres qu'elles concernent --- la TVA, l'impôt sur les sociétés, le travail.

% ===========================================================================
\section{Ce que font les autres pays}

Quatre pays ont mené depuis la \emph{Mirrlees Review} un exercice comparable, et leurs recommandations se recoupent largement. Aucun de ces
rapports n'a été écrit par les mêmes personnes, aucun n'a le même mandat, et presque tous
recommandent d'élargir les assiettes, de réduire le nombre de taux, de déplacer la charge du travail
vers la consommation et le patrimoine, et de taxer le foncier sur sa valeur réelle
(figure~\ref{fig:convergence} et tableau~\ref{tab:revues}).

\figbloc{c20_convergence_revues}{Cinq revues indépendantes, et des recommandations qui se
recoupent largement}{fig:convergence}

\begin{table}[htbp]
\caption{Cinq revues fiscales d'ensemble et leurs recommandations principales}
\label{tab:revues}
\small
\begin{tabular}{@{}p{3.1cm}p{11.2cm}@{}}
\toprule
\textbf{Royaume-Uni}\newline 2010-2011 & \emph{Mirrlees Review}, Institute for Fiscal Studies.
TVA à base large et taux unique, compensée par le barème direct. Fusion de l'impôt sur le revenu
et des cotisations. Exonération du rendement normal de l'épargne, allocation pour fonds propres
des sociétés. Impôt foncier assis sur la valeur, en remplacement des droits de mutation. Prix
unique du carbone.\\
\midrule
\textbf{Nouvelle-Zélande}\newline 2018-2019 & \emph{Tax Working Group}. Maintien d'une TVA à base
large et taux uniforme, considérée comme un acquis à ne pas abîmer. Le mandat excluait
explicitement toute hausse du taux.\\
\midrule
\textbf{Irlande}\newline 2022 & \emph{Commission on Taxation and Welfare}, «~Foundations for the
Future~». 116 recommandations. Déplacer la charge du travail vers le capital, le patrimoine et la
consommation. Fusionner les taux réduits de TVA et les relever progressivement. Imposer la valeur
du terrain et relever la taxe foncière locale. Supprimer les effets de seuil des prestations.\\
\midrule
\textbf{Norvège}\newline 2022 & Comité Torvik, NOU 2022:20. Conserver l'impôt sur la fortune mais
baisser son taux, élargir son assiette et supprimer toutes les décotes d'évaluation pour valoriser
les actifs à 100~\% de leur valeur de marché. Rétablir un impôt sur les successions. Réduire de
trois à deux le nombre de taux de TVA, voire aller au taux unique de 25~\%.\\
\midrule
\textbf{Pays-Bas}\newline 2024 & «~Belastingen in maatschappelijk perspectief~», groupe de travail
interministériel. Constat que l'instrumentalisation fiscale est allée trop loin. Faire contribuer
davantage les retraités en fiscalisant les cotisations de retraite de base. Éteindre la subvention
au logement occupé par son propriétaire. Imposer plus uniformément les revenus du capital et les
transmissions. Mieux tarifer les externalités. Le produit sert à alléger la charge pesant sur le
travail.\\
\bottomrule
\end{tabular}
\end{table}

Cette convergence ne prouve pas que les recommandations soient justes : cinq groupes
d'économistes formés dans la même tradition peuvent se tromper ensemble. Elle établit en revanche
qu'il ne s'agit pas de préférences nationales ni de positionnements politiques, puisque les mandats,
les majorités commanditaires et les systèmes de départ différaient. Ce qui converge, ce sont les
conclusions tirées d'un même corps de résultats appliqué à des situations différentes.

Deux de ces rapports posent frontalement une question que la France n'ouvre pas : la
contribution des retraités. Le rapport néerlandais recommande de fiscaliser les cotisations de
retraite de base, et la commission irlandaise de rééquilibrer entre générations. En France, le
niveau de vie médian des 65 ans et plus atteint 94~\% de celui de l'ensemble de la population,
contre 90~\% dans l'Union, 83~\% en Allemagne et 78~\% au Danemark ; la pauvreté frappe 21,4~\% des
moins de 18 ans contre 12,4~\% des plus de 65 ans, soit l'inverse du profil allemand
(figure~\ref{fig:retraites}).

\figbloc{p1_retraites_niveau_vie}{Niveau de vie relatif des 65 ans et plus, et taux de pauvreté par
âge, 2024}{fig:retraites}

Ce constat ne dicte aucune conclusion. Le niveau de vie relatif des retraités résulte de
décisions passées, prises en connaissance de cause, sur lesquelles les intéressés ont fondé des
plans de long terme --- et le respect des attentes légitimes figure parmi les critères d'équité de la
\emph{Mirrlees Review}. Le constat impose seulement que la question soit traitée comme une question
de répartition entre générations, au lieu d'être évacuée.

\subsection{Ce que mesure un indice}

L'indice de compétitivité fiscale de la Tax Foundation classe la France au dernier rang des
trente-huit pays de l'OCDE en 2025, avec un score de 45,8 sur 100. Le détail est plus instructif
que le rang global : 38\textsuperscript{e} sur l'impôt des sociétés, 34\textsuperscript{e} sur les
impôts fonciers, 31\textsuperscript{e} sur l'impôt sur le revenu et sur les taxes à la consommation,
mais 12\textsuperscript{e} sur les règles fiscales internationales (figure~\ref{fig:indice}).

\figbloc{c9_matrice_indice}{La France dernière de l'OCDE, mais pas sur tous les
tableaux}{fig:indice}

\begin{encadre}{Un indice est un argument déguisé en mesure}
Celui-ci agrège plus de quarante variables selon une pondération que ses auteurs ont choisie, et
cette pondération encode une préférence : pour la neutralité à l'égard de l'investissement et pour
une charge faible sur le capital. La redistribution n'y figure pas comme objectif, ce qui n'est pas
un oubli mais une position.

D'où la façon de le lire. Quand l'indice pénalise les impôts de production, il mesure une
distorsion que le théorème d'efficience productive établit indépendamment de toute préférence
politique. Quand il pénalise la progressivité de l'impôt sur le revenu ou l'existence d'un impôt
sur le patrimoine, il applique un jugement que la théorie ne tranche pas. Le premier usage est
légitime, le second doit être discuté comme ce qu'il est.
\end{encadre}

L'apport le plus utile de cette institution tient à ses données et à ses façons de les montrer. Plusieurs figures de ce cours en viennent : la cascade des cent
euros de profit, le pyramidage d'une taxe sur le chiffre d'affaires, la valeur actuelle des
amortissements, les taux d'impôt sur les sociétés depuis 1980, l'assiette de la TVA. On peut
emprunter l'instrument sans souscrire à l'usage qui en est fait.

% ===========================================================================
\section{Ce qu'il faudrait changer}

Le tableau~\ref{tab:synthese} reprend la forme du tableau final de la \emph{Mirrlees Review},
qui mettait en regard un bon système et le système britannique ; la colonne du milieu y remplace
l'idéal par le chiffre qui situe la France.

\begin{table}[htbp]
\caption{Impôt par impôt : pourquoi il coûte, où en est la France, ce qu'il faudrait faire}
\label{tab:synthese}
\footnotesize
\renewcommand{\arraystretch}{1.3}
\begin{tabularx}{\linewidth}{@{}>{\raggedright\arraybackslash}p{2.2cm}
  >{\raggedright\arraybackslash}X>{\raggedright\arraybackslash}X>{\raggedright\arraybackslash}X@{}}
\toprule
\textbf{Impôt} & \textbf{Pourquoi il coûte} & \textbf{La France, en un chiffre} &
\textbf{Ce qu'il faudrait faire}\\
\midrule
Production & Il frappe ce qui sert à produire ; assis sur le chiffre d'affaires, il se cumule le
long de la chaîne & 77~Md\euro{} sans fonction corrective ; taxes sur la masse salariale à
1,9~\% du PIB, troisième de l'OCDE & Supprimer, en commençant par le chiffre d'affaires ;
financer par des assiettes larges\\
\midrule
Sociétés & Il s'ajoute à l'impôt sur le dividende, taxe le rendement normal et favorise la dette &
Taux de 25,8~\% pour la plupart des entreprises, 36,1~\% pour les plus grandes en 2025 et 2026 ;
recette de 2,35~\% du PIB, septième plus basse &
Taux stable, sans surtaxe ; déduction d'un rendement normal sur les fonds propres nouveaux\\
\midrule
Travail & Le taux sur l'euro suivant décide, et l'empilement des dispositifs le rend le plus lourd en
bas & Coin marginal de 64,6~\% à 67~\% du salaire moyen ; taux supérieur appliqué à partir de treize
fois ce salaire & Lisser le creux au-dessus du salaire minimum ; fusionner impôt sur le revenu et
CSG\\
\midrule
TVA & Chaque taux réduit est une subvention qui donne plus d'euros aux ménages aisés & Taux normal
de 20~\%, dans la moyenne ; assiette effective de 51~\%, contre 55~\% & Rapprocher les taux réduits
du taux normal, avec un transfert forfaitaire\\
\midrule
Épargne & Taxer le rendement taxe l'attente, et l'inflation en plus & De 0 à 97~\% selon le
support ; 34~\% sur dividendes et plus-values, contre 25 et 20~\% en moyenne & Abattement pour
rendement normal, assiette indexée, un seul régime\\
\midrule
Immobilier et transmissions & Le sol se taxe sans dommage, la transaction pénalise la mobilité ;
des taux élevés sur une assiette étroite incitent à l'optimisation & Détention à 1,9~\% et
transactions à 0,7~\% du PIB, le double de la moyenne ; successions à 0,76~\%, premier rang de
l'OCDE & Réviser les valeurs locatives ; déplacer la charge des droits de mutation vers le sol ;
élargir l'assiette des successions et baisser les taux\\
\midrule
Carbone & Des prix différents rendent plus chère chaque tonne évitée & Plus de 120~\euro{} la tonne
sur la route ; 41~\% des émissions de l'industrie sans aucun prix & Étendre le prix sans le relever ;
quotas aux enchères ; reversement forfaitaire\\
\bottomrule
\end{tabularx}
\end{table}

\subsection{Dix écarts, rangés par certitude du gain}

\begin{enumerate}[leftmargin=1.5em,itemsep=0.35em]
\item \textbf{Les 77~Md\euro{} assis sur des facteurs de production.} Le seul gain qui ne dépend
d'aucune hypothèse de comportement.
\item \textbf{La dispersion du prix du carbone.} Même réduction d'émissions, coût inférieur, sans
accord préalable sur l'ambition climatique.
\item \textbf{Les valeurs locatives de 1970.} Préalable technique à toute réforme du foncier : tant
que l'assiette est fausse, toute réforme amplifierait des écarts injustifiés.
\item \textbf{Les droits sur les transactions, au double de la moyenne de l'OCDE.} Ils découragent les
déménagements et suivent le marché ; leur produit peut être repris par l'impôt récurrent.
\item \textbf{Le creux de taux marginal au-dessus du salaire minimum.} Personne ne l'a voulu ; il
résulte de l'empilement.
\item \textbf{Les taux réduits de TVA.} Réforme la mieux documentée et la plus rejetée ; elle n'a de
sens qu'accompagnée de sa compensation chiffrée.
\item \textbf{Le rendement normal et l'inflation dans la fiscalité de l'épargne.} Le même rendement
réel supporte de 0 à 97~\% selon l'enveloppe.
\item \textbf{Deux impôts sur le revenu qui ne partagent ni assiette, ni unité, ni barème.}
153~Md\euro{} de contribution sociale généralisée et 96~Md\euro{} d'impôt sur le revenu.
\item \textbf{L'impôt sur les sociétés.} Biais en faveur de la dette, taux instable, et un écart de
trente points avec le taux supérieur de l'impôt sur le revenu qui pousse à loger les revenus dans
des sociétés.
\item \textbf{Les transmissions.} Le rendement le plus élevé de l'OCDE, obtenu avec des taux élevés
sur une assiette étroite.
\end{enumerate}

Les quatre premiers écarts reposent sur des mécanismes qui ne dépendent d'aucune estimation : supprimer une taxe sur un
intrant améliore l'allocation, un prix unique du carbone atteint un objectif donné à moindre coût,
une assiette fausse produit des écarts arbitraires, une taxe sur la transaction décourage les
déménagements. Les suivants supposent des réactions de comportement dont l'ampleur se discute. Les deux
derniers engagent des arbitrages où des économistes compétents divergent.

\subsection{Ce qu'il ne faut pas changer}

Le niveau global de prélèvement n'est pas le premier problème. Le Danemark prélève
davantage avec une structure presque opposée. Si le niveau était la variable décisive, les deux pays
devraient présenter les mêmes difficultés.

Le taux normal de TVA n'a pas besoin d'être relevé. Il est dans la moyenne de l'OCDE ; c'est
l'assiette qui est plus étroite que la moyenne, et c'est elle qu'il faut élargir.

Le prix du carbone sur les carburants routiers n'a pas besoin d'être relevé ; il a besoin
d'être étendu. C'est le seul endroit où la France applique déjà un prix élevé et cohérent avec ses
objectifs.

Les règles d'amortissement de l'impôt sur les sociétés ne sont pas le problème. La valeur
actuelle des déductions est supérieure à la moyenne de l'OCDE pour les trois catégories d'actifs.

Le taux supérieur de l'impôt sur le revenu n'est pas le premier problème du travail. Il ne
s'applique qu'à partir de treize fois le salaire moyen ; le prélèvement qui pèse sur le plus grand
nombre est celui du bas de l'échelle.

\subsection{Les principes qui guident ces propositions}

Ces propositions séparent ce qui est démontré de ce qui est supposé. Elles reposent sur quatre
sortes de fondements : un calcul qu'on peut refaire au crayon, comme la perte qui croît avec le
carré du taux ou la cascade d'une taxe sur le chiffre d'affaires ; un fait mesuré, comme la mobilité
des revenus ou le coût des niches ; un effet causal identifié, comme la part de l'impôt sur les
sociétés qui retombe sur les salaires ou l'effet des droits de mutation sur les déménagements ;
enfin des hypothèses plausibles que personne n'a encore testées. Le classement des dix écarts va
du plus démontré au moins démontré. Le cours formule cinq de ces hypothèses : que la C3S freine les
effets d'agglomération, que la subvention au propriétaire occupant coûte de la croissance, que
l'imposition commune encourage la formation des couples, qu'un taux plus bas sur une assiette plus
large rapporte davantage, et qu'un impôt dû même sans bénéfice pousse à monter en gamme --- cette
dernière se heurtant au fait qu'une taxe sur le chiffre d'affaires frappe la marge et non la
productivité. Aucune n'est démontrée ; chacune mériterait une évaluation sur données françaises
plutôt qu'un emprunt aux études étrangères.

Elles se jugent sur le cas courant et sur l'ordre de grandeur. Un taux d'impôt sur les sociétés de
36,1~\% existe, mais pour environ 450 entreprises ; le cas courant est 25~\%. Une perte de 360 à
720~millions d'euros paraît faible, mais elle représente 10 à 20 centimes par euro prélevé par la
C3S. Un argument fiscal doit dire qui il concerne et combien il pèse, rapporté à ce que l'impôt
rapporte.

Un résultat théorique ne s'applique que si ses hypothèses tiennent. Le théorème d'efficience
productive suppose que l'État taxe en aval tout ce qu'il veut taxer ; il ne sait atteindre ni les
bénéfices que leurs propriétaires laissent dans une société, ni ceux des actionnaires étrangers, et
c'est ce qui justifie l'impôt sur les sociétés. L'impôt forfaitaire est le plus efficient des impôts
dans le modèle ; en Angleterre, il a duré trois ans.

Il faut aussi regarder comment on échappe à l'impôt. Travailler une heure de moins est une réaction
que l'assiette la mieux construite ne supprime pas. Requalifier un revenu, déplacer un bénéfice,
diviser une société pour rester sous un abattement en sont d'autres, qui dépendent de ce que la loi
permet et qu'une assiette sans issue réduit.

Chaque proposition vient avec son coût et son financement. La déduction pour fonds propres améliore
l'impôt sur les sociétés en principe ; la Belgique a montré ce qu'elle coûte quand elle porte sur
tout le stock de fonds propres. Remplacer la CVAE par l'impôt sur les sociétés réduit une distorsion
mais reporte la charge sur une base plus mobile. L'argent public étant fongible, l'affectation d'une
recette à une politique n'est pas une raison de garder une mauvaise assiette ; elle oblige seulement
à dire ce qui la remplace.

Reste la règle la plus difficile à défendre : la progressivité se juge sur l'ensemble des
prélèvements et des transferts. Une hausse de TVA compensée par un transfert peut améliorer la
situation des plus modestes ; jugée seule, elle sera rejetée. Toute réforme de ce cours doit donc
être présentée avec sa compensation, chiffrée par dixième de niveau de vie, faute de quoi elle sera
jugée sur la moitié de ses effets.

\subsection{Quand la théorie et la pratique divergent, qui doit apprendre ?}

Mankiw, Weinzierl et Yagan concluent sur un bilan en deux colonnes. Certaines évolutions
ressemblent à des victoires de la théorie : la généralisation de la TVA, la baisse des taux
statutaires sur le capital, l'aplatissement des barèmes. D'autres recommandations n'apparaissent
nulle part et n'y apparaîtront pas : l'impôt dépendant de caractéristiques innées, l'impôt sur le
capital variant avec les accidents du revenu du travail. «~Peu d'économistes conseillant un candidat
ou un ministre auraient l'audace d'avancer ces idées~» (figure~\ref{fig:lecons}).

\figbloc{c13_lecons_mankiw}{Les huit leçons de la théorie, et ce que la pratique en a
fait}{fig:lecons}

Deux explications sont possibles ; les auteurs jugent la seconde au moins aussi plausible que la première. Soit la théorie a raison, et le public comme les gouvernements tardent à
comprendre des résultats contre-intuitifs. Soit la tradition des finances publiques contient des
idées que la théorie moderne a laissées de côté : le principe du bénéfice et l'équité horizontale,
qu'aucune fonction de bien-être standard ne contient.

Là où la théorie établit un mécanisme
qu'on peut refaire au crayon --- la perte sèche qui croît avec le carré du taux, la distorsion
produite par un impôt sur un intrant, l'arithmétique qui convertit un impôt sur le stock en taux sur
le flux, le fait qu'un taux réduit donne plus d'euros aux riches --- elle l'emporte sur l'habitude,
et l'écart avec la pratique est un défaut à corriger. Là où elle produit une recommandation que
toutes les sociétés refusent, l'hypothèse de travail est qu'elle omet quelque chose. Le critère
sépare les résultats qui tiennent par un calcul de ceux qui tiennent par une fonction de bien-être
qu'on a choisie : les premiers sont contraignants pour tout le monde, les seconds ne le sont que pour
qui accepte la fonction, et c'est exactement ce dont on débat en politique.

% ===========================================================================
\vspace{1.4em}
{\color{bleu}\hrule height 0.8pt}
\vspace{0.8em}

{\small
\textbf{Les trois textes.} N. Gregory Mankiw, Matthew Weinzierl et Danny Yagan, «~Optimal Taxation
in Theory and Practice~», \emph{Journal of Economic Perspectives} 23(4), 2009, p.~147-174
\textbullet{} James Mirrlees \emph{et al.}, \emph{Tax by Design}, Institute for Fiscal Studies et
Oxford University Press, 2011 \textbullet{} Åsa Johansson, Christopher Heady, Jens Arnold, Bert Brys
et Laura Vartia, «~Taxation and Economic Growth~», OCDE, document de travail du département des
affaires économiques n°~620, 2008, et la mise en cause de sa robustesse par Jing Xing, «~Tax
structure and growth: how robust is the empirical evidence~?~», \emph{Economics Letters}, 2012.

\textbf{Résultats cités nommément.} James Mirrlees (1971) \textbullet{} Peter Diamond et James
Mirrlees (1971) \textbullet{} Anthony Atkinson et Joseph Stiglitz (1976) \textbullet{} George Akerlof
(1978) \textbullet{} Matti Tuomala (1990), Peter Diamond (1998) et Emmanuel Saez (2001, 2002)
\textbullet{} Christophe Chamley (1986) et Kenneth Judd (1985), Ludwig Straub et Iván Werning (2020)
\textbullet{} Narayana Kocherlakota (2005) \textbullet{} David Carey et Josette Rabesona, «~Tax
Ratios on Labour and Capital Income and on Consumption~», \emph{OECD Economic Studies} n°~35, 2002,
repris dans P.~B. Sørensen (dir.), \emph{Measuring the Tax Burden on Capital and Labor}, MIT Press,
2004 \textbullet{} Youssef Benzarti et Dorian Carloni, «~Who Really Benefits from Consumption Tax
Cuts? Evidence from a Large VAT Reform in France~», \emph{American Economic Journal: Economic
Policy}, 2019.

\textbf{Autres travaux cités.} Jules Dupuit, «~De la mesure de l'utilité des travaux publics~»,
\emph{Annales des ponts et chaussées}, 1844 \textbullet{} Timothy Besley, Ian Preston et Michael
Ridge, «~Fiscal anarchy in the UK: Modelling poll tax noncompliance~», \emph{Journal of Public
Economics}, 1997 \textbullet{} Emmanuel Saez, Joel Slemrod et Seth Giertz, «~The Elasticity of
Taxable Income with Respect to Marginal Tax Rates: A Critical Review~», \emph{Journal of Economic
Literature}, 2012 \textbullet{} Joel Slemrod et Wojciech Kopczuk, «~The optimal elasticity of
taxable income~», \emph{Journal of Public Economics}, 2002 \textbullet{} Pierre-Yves Cabannes,
Cédric Houdré et Camille Landais, «~Comment le revenu imposable des ménages réagit-il à sa
taxation ?~», \emph{Économie et Statistique} n°~467-468, 2014 \textbullet{} Étienne Lehmann,
François Marical et Laurence Rioux, «~Labor income responds differently to income-tax and
payroll-tax reforms~», \emph{Journal of Public Economics}, 2013 \textbullet{} Thomas Loisel et
Michaël Sicsic, «~Peu de mobilité dans l'échelle des revenus entre 2003 et 2019~», \emph{Insee
Analyses} n°~82, 2023, et «~La mobilité des individus le long de l'échelle des revenus en France
sur la période 2003-2021~», \emph{Économie et Statistique} n°~545, 2024 \textbullet{} Camille
Landais, «~Le quotient familial a-t-il stimulé la natalité française ?~», \emph{Économie publique}
n°~13, 2003 \textbullet{} Mathias André, «~L'imposition conjointe des couples mariés et pacsés
organise une redistribution en direction des couples les plus aisés~», \emph{France, portrait
social}, Insee, 2019 \textbullet{} Conseil des prélèvements obligatoires, «~Conforter l'égalité des
citoyens devant l'imposition des revenus~», octobre 2024 \textbullet{} Clemens Fuest, Andreas Peichl et Sebastian Siegloch, «~Do Higher
Corporate Taxes Reduce Wages? Micro Evidence from Germany~», \emph{American Economic Review}, 2018
\textbullet{} Juan Carlos Suárez Serrato et Owen Zidar, «~Who Benefits from State Corporate Tax
Cuts?~», \emph{American Economic Review}, 2016 \textbullet{} Sebastian Beer, Ruud de Mooij et Li Liu,
«~International Corporate Tax Avoidance: A Review of the Channels, Magnitudes, and Blind Spots~»,
\emph{Journal of Economic Surveys}, 2020 \textbullet{} John Coles, Elena Patel, Nathan Seegert et
Matthew Smith, «~How Do Firms Respond to Corporate Taxes?~», \emph{Journal of Accounting Research},
2022 \textbullet{} Thomas Tørsløv, Ludvig Wier et Gabriel Zucman, «~The Missing Profits of
Nations~», \emph{Review of Economic Studies}, 2023 \textbullet{} Kimberly Clausing, «~Corporate tax
revenues in OECD countries~», \emph{International Tax and Public Finance}, 2007 \textbullet{} Laura
Kawano et Joel Slemrod, «~How do corporate tax bases change when corporate tax rates change?~»,
\emph{International Tax and Public Finance}, 2016 \textbullet{} Michael Devereux, Rachel Griffith
et Alexander Klemm, «~Corporate income tax reforms and international tax competition~»,
\emph{Economic Policy}, 2002 \textbullet{} Ruud de Mooij et Gaëtan Nicodème, «~Corporate tax policy
and incorporation in the EU~», \emph{International Tax and Public Finance}, 2008 \textbullet{} Ruud
de Mooij et Michael Devereux, «~An applied analysis of ACE and CBIT reforms in the EU~»,
\emph{International Tax and Public Finance}, 2011 \textbullet{} Ruud de Mooij, «~Tax Biases to Debt
Finance~», FMI, \emph{Staff Discussion Note} 11/11, 2011 \textbullet{} Muhammad Meki, «~Levelling
the debt-equity playing field: Evidence from Belgium~», \emph{European Economic Review}, 2023
\textbullet{} Shafik Hebous et Martin Ruf, «~Evaluating the effects of ACE systems on multinational
debt financing and investment~», \emph{Journal of Public Economics}, 2017 \textbullet{} Nicola
Branzoli et Antonella Caiumi, «~How effective is an incremental ACE in addressing the debt bias?~»,
\emph{International Tax and Public Finance}, 2020 \textbullet{} Benjamin Hansen, Keaton Miller et
Caroline Weber, «~Vertical integration and production inefficiency in the presence of a gross
receipts tax~», \emph{Journal of Public Economics}, 2022 \textbullet{} Fatih Guvenen \emph{et al.},
«~Use It or Lose It: Efficiency and Redistributional Effects of Wealth Taxation~», \emph{Quarterly
Journal of Economics}, 2023 \textbullet{} Fonds monétaire international, «~A Fair and Substantial
Contribution by the Financial Sector~», rapport au G20, 2010 \textbullet{} Martin Gervais, «~Housing
taxation and capital accumulation~», \emph{Journal of Monetary Economics}, 2002 \textbullet{}
Jukka-Pekka Laamanen, «~Home-ownership and the labour market: Evidence from rental housing market
deregulation~», \emph{Labour Economics}, 2017 \textbullet{} David Blanchflower et Andrew Oswald,
«~Does High Home-Ownership Impair the Labor Market?~», NBER, 2013 \textbullet{} Christian Hilber et
Teemu Lyytikäinen, «~Transfer taxes and household mobility~», \emph{Journal of Urban Economics},
2017 \textbullet{} Michael Best et Henrik Kleven, «~Housing Market Responses to Transaction Taxes~»,
\emph{Review of Economic Studies}, 2018 \textbullet{} Guillaume Bérard et Alain Trannoy, «~The
impact of the 2014 increase in the real estate transfer taxes on the French housing market~»,
\emph{Économie et Statistique}, 2018.

\textbf{Rapports publics français.} \emph{Conseil d'analyse économique} : Yannick L'Horty, Philippe
Martin et Thierry Mayer, «~Baisses de charges : stop ou encore ?~», note n°~49, janvier 2019
\textbullet{} Dominique Bureau, Fanny Henriet et Katheline Schubert, «~Pour le climat : une taxe
juste, pas juste une taxe~», note n°~50, mars 2019 \textbullet{} Philippe Martin et Alain Trannoy,
«~Les impôts sur (ou contre) la production~», note n°~53, juin 2019 \textbullet{} Clément
Dherbécourt, Gabrielle Fack, Camille Landais et Stefanie Stantcheva, «~Repenser l'héritage~», note
n°~69, décembre 2021. \emph{Conseil des prélèvements obligatoires} : «~La taxe sur la valeur ajoutée
(TVA), un impôt à recentrer sur son objectif de rendement pour les finances publiques~», février
2023 \textbullet{} «~Pour une fiscalité du logement plus cohérente~», décembre 2023
\textbullet{} «~Tracer un cadre fiscal et social pluriannuel pour l'industrie française~», septembre
2025. \emph{Cour des comptes} : «~Maîtriser la dynamique des allègements généraux de cotisations
sociales, contribuer à l'équilibre financier de la sécurité sociale~», chapitre~III du rapport
\emph{La sécurité sociale} sur l'application des lois de financement, mai 2025 \textbullet{} «~Le
pacte Dutreil : un dispositif fiscal en forte croissance à mieux cibler~», rapport public
thématique, novembre 2025 \textbullet{} \emph{Le budget de l'État en 2025 : résultats et gestion},
avril 2026.

\textbf{Données.} OCDE : Revenue Statistics (DSD\_REV\_COMP\_OECD), base de données fiscale (taux
supérieur et seuil de l'impôt sur le revenu), \emph{Taxing Wages}, modèle TaxBEN, \emph{Effective
Carbon Rates}, \emph{Corporate Tax Statistics} \textbullet{} Tax Foundation : \emph{Corporate Tax
Rates around the World} et \emph{International Tax Competitiveness Index 2025}, données publiées sur
github.com/TaxFoundation \textbullet{} Eurostat : \emph{National Tax Lists}, \texttt{gov\_10a\_taxag},
\texttt{nama\_10\_gdp}, \texttt{ilc\_pnp2}, \texttt{ilc\_li02} \textbullet{} rapports d'Irlande
(2022), de Norvège (NOU 2022:20), des Pays-Bas (2024) et de Nouvelle-Zélande (2019).

\textbf{Méthode.} Les figures construites par le calcul --- perte au carré, sommet des recettes,
pyramidage, amortissement, taux réduit, taux unique, taux supérieur, impôt sur l'attente, inflation,
supports d'épargne, stock et flux, capitalisation --- ne dépendent d'aucune estimation économétrique ;
leurs formules figurent sous chacune. Les autres reposent sur les sources ci-dessus ; le script
\texttt{extract\_reelles.py} télécharge les données de l'OCDE et de la Tax Foundation, et chaque
chiffre du cours renvoie à un fichier du dossier \texttt{donnees}.}

\end{document}
